% ============================================================
%  CHAPTER V -- Interpolation of Operators   (printed pp. 177-216)
%  transcribed from src/ch5.blocks.md (PDF pp. 189-228)
% ============================================================
\chapter{Interpolation of Operators}

Let $T$ be a linear operator mapping some linear space $\mathcal{A}$ into
another linear space $\mathcal{B}$. Suppose $A_{0}$, $A_{1}$ are Banach
subspaces of $\mathcal{A}$ and $B_{0}$, $B_{1}$ are Banach subspaces of
$\mathcal{B}$ such that when $T$ is restricted to $A_{i}$ it maps this space
continuously into $B_{i}$, $i = 0$, 1. When this is the case, it can be shown
often that there exist infinitely many pairs of ``intermediate'' Banach spaces
($A$, $B$), with $A \subset \mathcal{A}$ and $B \subset \mathcal{B}$, such that
$T$ restricted to $A$ is a bounded operator mapping into $B$. General theorems
of this kind are called interpolation theorems for linear operators and they
have many important applications in Fourier analysis. We have already
encountered an example of this situation when $\mathcal{A} = \mathcal{B}$ is
the space of tempered distributions, $T$ is the Fourier transform,
$A_{0} = L^{1}(E_{n})$, $B_{0} = L^{\infty}(E_{n})$ and
$A_{1} = L^{2}(E_{n}) = B_{1}$; it is then true that if $T$ is restricted to
$L^{p}(E_{n})$, $1 \leq p \leq 2$, it maps this space boundedly into
$L^{q}(E_{n})$, where $(1/p) + (1/q) = 1$. This follows from a basic
interpolation theorem concerning $L^{p}$ spaces: the $M$. Riesz interpolation
theorem. We shall study this theorem and some of its applications in the first
section of this chapter. There is another type of interpolation theorem in
which we do not need assumptions as strong as boundedness at the
``end-points'' $(A_{0}, B_{0})$ and $(A_{1}, B_{1})$; moreover, we can assume
certain conditions weaker than the linearity of $T$. This theorem, known as the
Marcinkiewicz theorem, applies to the maximal function operator $f \to m_{f}$
that was studied in the second chapter. The second section of this chapter is
devoted to applications of the Marcinkiewicz theorem. This theorem is best
understood when it is extended to a collection of function spaces that
includes the $L^{p}$ classes. This collection, the $L(p, q)$ spaces, as well as
the extended version of the Marcinkiewicz theorem is the subject of the third
section. In the fourth section we develop an extension of the M. Riesz theorem
in which the operators vary, in a sufficiently smooth manner, with the
``intermediate'' pairs of Banach spaces.

\section{The M. Riesz Convexity Theorem and Interpolation of Operators Defined
on \texorpdfstring{$L^{p}$}{Lp} Spaces}\index{interpolation@interpolation!of operators@of operators}

In the previous chapters we encountered several operators acting on $L^{p}$
spaces. One class of such operators consists of the convolution operators
determined by functions in $L^{p}(E_{n})$; that is, if we fix
$f \in L^{p}(E_{n})$ we obtain the operator $T$ that maps $g \in L^{1}(E_{n})$
into the function $f * g$. From Theorem~\ref{thm:c1-1-3} of Chapter I we see
that $T$ is a bounded operator defined on $L^{1}(E_{n})$ having values in
$L^{p}(E_{n})$; furthermore, the operator norm,
$\|T\|$,\footnote{If $T$ is a bounded transformation from $L^{r}$ into $L^{s}$
we shall denote its norm by $\|T\|$. That is,
$\|T\| = \inf\{k \geq 0 : \|Tf\|_{s} \leq k\|f\|_{r} \text{ for all } f \in L^{r}\}$.}
does not exceed $\|f\|_{p}$. On the other hand, an immediate application of
H\"older's inequality shows that $T$ also maps $L^{p'}(E_{n})$ boundedly into
$L^{\infty}(E_{n})$ when $(1/p) + (1/p') = 1$; as in the last case, the
operator norm $\|T\|$ does not exceed $\|f\|_{p}$. Since $T$ is defined on
$L^{1}(E_{n})$ and on $L^{p'}(E_{n})$ there is a natural extension of $T$ to
all the spaces $L^{r}(E_{n})$ when $1 \leq r \leq p'$: If $g \in L^{r}(E_{n})$
then we can find $g_{1} \in L^{1}(E_{n})$ and $g_{2} \in L^{p'}(E_{n})$ such
that $g = g_{1} + g_{2}$ and we can define $Tg$ to be $Tg_{1} + Tg_{2}$ (it is
clear that $Tg$ does not depend on the particular splitting of $g$ as
$g = g_{1} + g_{2}$).\footnote{For example, by letting $g_{1}(x) = g(x)$ when
$|g(x)| \geq 1$ and $g_{1}(x) = 0$ when $|g(x)| < 1$, we obtain such a
decomposition.} It is natural, therefore, to ask whether $T$ maps $L^{r}(E_{n})$
boundedly into some space $L^{q}(E_{n})$. We have already stated without proof
(see \eqref{eq:c1-4-3} in the first chapter) that this is the case when
$(1/q) = (1/p) + (1/r) - 1$, and $\|T\|$ is, again, less than or equal to
$\|f\|_{p}$. That is,
\begin{equation}\label{eq:c5-1-1}
\|f*g\|_{q}\leq\|f\|_{p}\|g\|_{r}
\end{equation}
whenever $f \in L^{p}(E_{n})$ and $g \in L^{r}(E_{n})$. We shall show, in this
section, that this inequality (known as Young's\index{Young inequality@Young inequality} inequality) follows easily from
a general result concerning linear operators acting on $L^{p}$-spaces.

This general result can also be applied to the Fourier transform
$\mathcal{F}$. We have treated $\mathcal{F}$ mostly as a linear operator
defined on $L^{1}(E_{n})$ and $L^{2}(E_{n})$. Theorem~\ref{thm:c1-1-1}, part
(a), of the first chapter asserts that the first space is mapped boundedly into
$L^{\infty}(E_{n})$ with operator norm $\leq 1$, while the Plancherel theorem
tells us that the second space is mapped onto itself with operator norm equal to
1. For the same reasons we gave in the previous paragraph, $\mathcal{F}$ has a
natural extension to the ``intermediate'' spaces $L^{p}(E_{n})$,
$1 \leq p \leq 2$ (see, also, the remarks made toward the end of Section 2 in
the first chapter). We shall show that this space is mapped by $\mathcal{F}$
boundedly into $L^{p'}(E_{n})$, where $(1/p) + (1/p') = 1$, with operator norm
$\leq 1$. That is, if $f \in L^{p}(E_{n})$, $1 \leq p \leq 2$, then
$\hat{f} = \mathcal{F}f$ belongs to $L^{p'}(E_{n})$ and
\begin{equation}\label{eq:c5-1-2}
\|\hat{f}\|_{p'}\leq\|f\|_{p}.
\end{equation}
This result is known as the inequality of Hausdorff-Young\index{Hausdorff-Young inequality@Hausdorff-Young inequality}.

In order to state the above mentioned general result concerning linear
operators, the $M$. Riesz convexity\index{Riesz convexity theorem@Riesz convexity theorem} theorem, we need to introduce some new
concepts and notation. Suppose $(M, \mathcal{M}, \mu)$ is a measure
space.\footnote{$M$ is a set, $\mathcal{M}$ the $\sigma$-algebra of measurable
subsets of $M$ and $\mu$ is a measure defined on $\mathcal{M}$. We shall
consider only $\sigma$-finite measure spaces; thus, all the standard results of
measure theory, such as the Radon-Nikodym theorem and the F. Riesz
representation theorem, are valid.} We let $L^{p}(M)$, $1 \leq p < \infty$,
denote the space of all those complex-valued functions $f$ satisfying
$\|f\|_{p} = (\int_{M}|f|^{p}\,d\mu)^{1/p} < \infty$. As is the case when
$M = E_{n}$ (see the beginning of Chapter I) we use the same notation,
$L^{p}(M)$, for the space of equivalence classes obtained by considering two
such functions to be equivalent when they agree almost everywhere; by defining
the norm of the class containing the function $f$ to be $\|f\|_{p}$ we then
obtain a Banach space. The same remarks apply to the space $L^{\infty}(M)$ of
essentially bounded measurable functions if we define $\|f\|_{\infty}$, for
$f \in L^{\infty}(M)$, to be the essential supremum of $f$.

If $f$ is a measurable function, a truncation\index{truncation@truncation} of $f$ is any one of the
functions $g$ defined by letting $g(x) = f(x)$ if $r_{1} < |f(x)| \leq r_{2}$
and $g(x) = 0$ otherwise, where $r_{1}$ and $r_{2}$ are nonnegative. Suppose
$T$ is an operator mapping a linear space $D$ of measurable functions on
$(M, \mathcal{M}, \mu)$ into measurable functions defined on another measure
space $(N, \mathcal{N}, \nu)$. We assume that $D$ contains all characteristic
functions of sets of finite measure and has the property that whenever
$f \in D$ and $g$ is a truncation of $f$, then $g \in D$. If there exists a
constant $k > 0$ such that
\[
\left(\int_{N}|Tf|^{q}\,d\nu\right)^{1/q}=\|Tf\|_{q}\leq k\|f\|_{p}
=k\left(\int_{M}|f|^{p}\,d\mu\right)^{1/p}
\]
for all $f \in D \cap L^{p}(M)$ we say that $T$ is of type $(p, q)$. The least
$k$ for which this inequality holds is called the $(p, q)$-norm of $T$.

The M. Riesz convexity theorem can now be formulated in the following way:

\begin{theorem}\label{thm:c5-1-3}
Suppose a linear operator $T$ is of type $(p_{i}, q_{i})$, with
$(p_{i}, q_{i})$-norm $k_{i}$, $i = 0$, 1, then $T$ is of type $(p_{t}, q_{t})$
with $(p_{t}, q_{t})$-norm $k_{t} \leq k_{0}^{1-t}k_{1}^{t}$ provided
$1/p_{t} = (1 - t)/p_{0} + t/p_{1}$ and $1/q_{t} = (1 - t)/q_{0} + t/q_{1}$,
with $0 \leq t \leq 1$.
\end{theorem}

Before proving this theorem we show how the two inequalities, \eqref{eq:c5-1-1}
and \eqref{eq:c5-1-2}, follow immediately from it. We observed that the
operator $T : g \to f * g$ is of types $(p_{0}, q_{0}) = (1, p)$ and
$(p_{1}, q_{1}) = (p', \infty)$ (its domain $D$ includes all the spaces
$L^{r}(E_{n})$, $1 \leq r \leq p'$); moreover, the $(1, p)$-norm and the
$(p', \infty)$-norm do not exceed $\|f\|_{p}$. Thus, $T$ is of type
$(p_{t}, q_{t})$, $0 \leq t \leq 1$, where
\[
\frac{1}{p_{t}}=\frac{1-t}{p_{0}}+\frac{t}{p_{1}}=1-t+\frac{t}{p'}=1-\frac{t}{p}
\]
and
\[
\frac{1}{q_{t}}=\frac{1-t}{q_{0}}+\frac{t}{q_{1}}=\frac{1-t}{p}+\frac{t}{\infty}
=\frac{1-t}{p}.%
\footnote{In case $p_{0}$, $p_{1}$, $q_{0}$ or $q_{1}$ is $\infty$ we adopt the
usual convention: We let $1/\infty = 0$.}
\]

Putting $r = p_{t}$ and $q = q_{t}$ we then have
\[
\frac{1}{q}=\frac{1-t}{p}=\frac{1}{p}+\left(1-\frac{t}{p}\right)-1
=\frac{1}{p}+\frac{1}{r}-1.
\]
Moreover, the $(r, q)$-norm of $T$ does not exceed
$\|f\|_{p}^{1-t}\|f\|_{p}^{t} = \|f\|_{p}$. But this is precisely inequality
\eqref{eq:c5-1-1}. Similarly, since the Fourier transform is of types
$(1, \infty)$ and $(2, 2)$, with norms not larger than 1, an application of
Theorem~\ref{thm:c5-1-3} with $t = 2/p'$ gives us the Hausdorff-Young inequality
\eqref{eq:c5-1-2}.

It is often useful to consider the following geometrical interpretation of the
M. Riesz convexity theorem: Let $Q$ be the square in the plane whose vertices
are the points $(0, 0)$, $(1, 0)$, $(1, 1)$, and $(0, 1)$. If $T$ is an operator
satisfying the hypotheses of Theorem~\ref{thm:c5-1-3}, we can associate with
$T$ the points $x_{0} = (1/p_{0}, 1/q_{0})$ and $x_{1} = (1/p_{1}, 1/q_{1})$ in
order to indicate that this operator is of types $(p_{0}, q_{0})$ and
$(p_{1}, q_{1})$. As $t$ ranges from 0 to 1, the point
$x_{t} = (1/p_{t}, 1/q_{t})$ traverses the line segment joining $x_{0}$ and
$x_{1}$; Theorem~\ref{thm:c5-1-3} asserts that, for each point
$(\alpha, \beta)$ of this segment, $T$ is of type $(1/\alpha, 1/\beta)$.

The justification for calling Theorem~\ref{thm:c5-1-3} a convexity theorem is
that, if $\varphi(t)$ is the logarithm of the $(p_{t}, q_{t})$-norm of $T$,
then $\varphi$ is a convex function. Moreover, the set of all points
$(1/p, 1/q)$ such that $T$ is of type $(p, q)$ is a convex set. We shall show
that Theorem~\ref{thm:c5-1-3} can be derived from a classical result in the
theory of functions whose conclusion also involves a logarithmically convex
function. This result, a Phragm\'en-Lindel\"of theorem (see \eqref{frp:c2-5-2}
of Chapter II), is known as the three lines\index{three lines theorem@three lines theorem} theorem and can be formulated in
the following way:

\begin{lemma}\label{lem:c5-1-4}
Suppose $F$ is a bounded continuous complex-valued function on the closed strip
$S = \{x + iy = z \in \mathbf{C} : 0 \leq x \leq 1\}$ that is analytic in the
interior of $S$. If $|F(iy)| \leq m_{0}$ and $|F(1 + iy)| \leq m_{1}$ for all
$y$ then $|F(x + iy)| \leq m_{0}^{1-x}m_{1}^{x}$ for all $z = x + iy \in S$.
That is, if we define
$k_{x} = \sup\{|F(x + iy)| : -\infty < y < \infty\}$ for $0 \leq x \leq 1$, then
$\varphi(x) = \log k_{x}$ is a convex function.
\end{lemma}

\begin{proof}
We can obviously assume that $m_{0}$ and $m_{1}$ are positive. Then by
considering the function whose values are $F(z)/m_{0}^{1-z}m_{1}^{z}$, if
necessary, we can reduce the problem to the case $m_{0} = 1 = m_{1}$. Thus, we
can suppose that $|F(iy)| \leq 1$ and $|F(1 + iy)| \leq 1$ for all $y$. We want
to show that $|F(z)| \leq 1$ for all $z \in S$. If we assume that
$\lim_{|y| \to \infty}F(x + iy) = 0$ uniformly for $0 \leq x \leq 1$ this
result follows easily from the maximum principle: for we can then find
$y_{0} > 0$ such that $|F(x + iy)| \leq 1$ for $|y| \geq y_{0}$, while
$|F(z)| \leq 1$ on the boundary of the rectangle with vertices $iy_{0}$,
$1 + iy_{0}$, $1 - iy_{0}$, $-iy_{0}$. In general, we can apply this result to
the function $F_{n}(z) = F(z)e^{(z^{2}-1)/n}$, $n = 1, 2, \ldots$, and, because
$F$ is bounded,
\[
\begin{aligned}&|F_{n}(z)|=|F(x+iy)|e^{-y^{2}/n}e^{(x^{2}-1)/n}\\&\qquad\leq|F(x+iy)|e^{-y^{2}/n}\to0
\quad\text{as}\quad y\to\infty,\end{aligned}
\]
uniformly for $0 \leq x \leq 1$, $|F_{n}(iy)| \leq 1$ and
$|F_{n}(1 + iy)| \leq 1$. Therefore, $|F_{n}(z)| \leq 1$ and letting $n$ tend
to $\infty$, we obtain the desired inequality.
\end{proof}

We now pass to the proof of Theorem~\ref{thm:c5-1-3}. In order to do this we
let $\alpha_{j} = 1/p_{j}$, $\beta_{j} = 1/q_{j}$, $j = 0$, 1,
$\alpha = 1/p_{t}$ and $\beta = 1/q_{t}$. Consequently, if
$\alpha(z) = (1 - z)\alpha_{0} + z\alpha_{1}$ and
$\beta(z) = (1 - z)\beta_{0} + z\beta_{1}$, for $z \in \mathbb{C}$, we have
$\alpha(j) = \alpha_{j}$, $\beta(j) = \beta_{j}$, $j = 0$, 1, $\alpha(t) = \alpha$
and $\beta(t) = \beta$.

First, we estimate $\|Tf\|_{q}$ for $f$ a finite linear combination of
characteristic functions of subsets of $M$ having finite measure. Such a simple
function, $f$, belongs to the domain of $T$ by definition. Since
\[
\|Tf\|_{q}=\sup\left|\int_{N}(Tf)g\,d\nu\right|,
\]
where the supremum is taken over all simple functions $g$ satisfying
$\|g\|_{q'} = 1$, $(1/q) + (1/q') = 1$, it suffices to show that the absolute
value of each such integral
\[
I=\int_{N}(Tf)g\,d\nu
\]
is less than or equal to $k_{0}^{1-t}k_{1}^{t}\|f\|_{p}$. Dividing by
$\|f\|_{p}$, if this norm is not zero, we can reduce the problem further to the
case $\|f\|_{p} = 1$.

Suppose, then, that $f = \sum_{j=1}^{m}a_{j}\chi_{E_{j}}$ and
$g = \sum_{k=1}^{n}b_{k}\chi_{F_{k}}$ are simple functions satisfying the above
conditions. We also assume that $p_{t} < \infty$ and $q_{t} > 1$; thus,
$\alpha > 0$ and $\beta < 1$. If $a_{j} = |a_{j}|e^{i\theta_{j}}$ and
$b_{k} = |b_{k}|e^{i\varphi_{k}}$ we define
\[
f_{z}=\sum_{j=1}^{m}|a_{j}|^{\alpha(z)/\alpha}e^{i\theta_{j}}\chi_{E_{j}}
\]
and
\[
g_{z}=\sum_{k=1}^{n}|b_{k}|^{(1-\beta(z))/(1-\beta)}e^{i\varphi_{k}}\chi_{F_{k}}
\qquad\text{for } z \in \mathbb{C}.
\]

One then obtains an entire function by putting
$F(z) = \int_{N}(Tf_{z})g_{z}\,d\nu$, and it follows from the definitions of
$\alpha(z)$ and $\beta(z)$ that $F(t) = I$. From the linearity of $T$ we see
that
\[
F(z)=\sum_{j,k=1}^{m,n}|a_{j}|^{\alpha(z)/\alpha}|b_{k}|^{(1-\beta(z))/(1-\beta)}\gamma_{jk},
\]
where
$\gamma_{jk} = e^{i(\theta_{j} + \varphi_{k})}\int_{N}(T\chi_{E_{j}})\chi_{F_{k}}\,d\nu$.
Since each of these summands is bounded in the strip $S$, it follows that the
entire function $F$ is bounded when restricted to this strip. If we can show
that $|F(iy)| \leq k_{0}$ and $|F(1 + iy)| \leq k_{1}$ for all $y$, then the
desired inequality, $|I| \leq k_{0}^{1-t}k_{1}^{t}$, will be an immediate
consequence of Lemma~\ref{lem:c5-1-4}.

We can obtain these two estimates by first observing that, since
$\alpha(iy) = \alpha_{0} + iy(\alpha_{1} - \alpha_{0})$ and
$1 - \beta(iy) = (1 - \beta_{0}) - iy(\beta_{1} - \beta_{0})$, we have
$|f_{iy}|^{p_{0}} = \bigl|e^{i\arg f}|f|^{iy(\alpha_{1} - \alpha_{0})/\alpha}|f|^{(p/p_{0})}\bigr|^{p_{0}} = |f|^{p}$
and
\[
|g_{iy}|^{q_{0}'}=\bigl|e^{i\arg g}|g|^{-iy(\beta_{1}-\beta_{0})/(1-\beta)}|g|^{(q'/q_{0}')}\bigr|^{q_{0}'}=|g|^{q'}.
\]
Thus, applying H\"older's inequality and the fact that $T$ is of type
$(p_{0}, q_{0})$ with norm $k_{0}$,
\begin{align*}
|F(iy)|&\leq\|Tf_{iy}\|_{q_{0}}\|g_{iy}\|_{q_{0}'}
\leq k_{0}\|f_{iy}\|_{p_{0}}\|g_{iy}\|_{q_{0}'}\\
&=k_{0}\left(\int_{M}|f|^{p}\,d\mu\right)^{1/p_{0}}
\left(\int_{N}|g|^{q'}\,d\nu\right)^{1/q_{0}'}\\
&=k_{0}\|f\|_{p}^{(p/p_{0})}\|g\|_{q'}^{(q'/q_{0}')}=k_{0}.
\end{align*}
A similar calculation shows that $|F(1 + iy)| \leq k_{1}$. We can conclude,
therefore, that $\|Tf\|_{q} \leq k_{0}^{1-t}k_{1}^{t}\|f\|_{p}$ for all simple
$f$ in $L^{p}(M)$.

In order to establish this inequality for the general
$f \in D \cap L^{p}(M)$, where $D$ is the domain of $T$, we shall show that we
can find a sequence $\{f_{n}\}$ of simple functions such that
$\|f_{n} - f\|_{p} \to 0$ and $(Tf_{n})(x) \to (Tf)(x)$ a.e. as $n \to \infty$.
If this were the case, an application of Fatou's lemma and the result we just
obtained for simple functions would give us the inequality
$\|Tf\|_{q} \leq \lim_{n \to \infty}\|Tf_{n}\|_{q} \leq \lim_{n \to \infty}\{k_{0}^{1-t}k_{1}^{t}\|f_{n}\|_{p}\} = k_{0}^{1-t}k_{1}^{t}\|f\|_{p}$,
and the theorem would be proved. In order to find such a sequence we can assume
$f \geq 0$ (for we may consider, separately, the positive and negative parts of
$\operatorname{Re}\{f\}$ and $\Im\{f\}$.) After relabelling, if necessary, we
can also assume that $p_{0} \leq p_{1}$. Let $f^{0}$ and $f^{1}$ be the
truncations of $f$ defined by letting
\[
f^{0}(x)=\begin{cases}
f(x) & \quad\text{when } f(x) > 1,\\
0 & \quad\text{when } f(x) \leq 1,
\end{cases}
\]
and $f^{1} = f - f^{0}$. Since $(f^{0})^{p_{0}} \leq f^{p}$,
$(f^{1})^{p_{1}} \leq f^{p}$ and $D$ contains all the truncations of $f$ we
have $f^{0} \in D \cap L^{p_{0}}(M)$ and $f^{1} \in D \cap L^{p_{1}}(M)$. If
$\{g_{m}\}$ is an increasing sequence of nonnegative simple functions
converging to $f$ it follows from the monotone convergence theorem of Lebesgue
that $\lim_{m \to \infty}\|g_{m} - f\|_{p} = 0$. For the same reason
$\|g_{m}^{0} - f^{0}\|_{p_{0}}$ and $\|g_{m}^{1} - f^{1}\|_{p_{1}}$ tend to 0
as $m$ tends to $\infty$, where $g_{m}^{0}$ and $g_{m}^{1}$ are the truncations
(at 1) obtained from $g_{m}$ in the same way as were $f^{0}$ and $f^{1}$
obtained from $f$. Since $T$ is of types $(p_{0}, q_{0})$ and $(p_{1}, q_{1})$
we must have $\|Tg_{m}^{0} - Tf^{0}\|_{q_{0}} \to 0$ and
$\|Tg_{m}^{1} - Tf^{1}\|_{q_{1}} \to 0$ as $m \to \infty$. Thus, there exists a
subsequence of $\{Tg_{m}^{0}\}$ converging almost everywhere to $Tf^{0}$.
Considering only those indices involved in this subsequence we can then find a
subsequence $\{g_{m}^{1}\}$ such that $Tg_{m}^{1}$ converges almost everywhere
to $Tf^{1}$. If we now let $f_{n} = g_{m}^{0} + g_{m}^{1}$ we have a sequence
$\{f_{n}\}$ satisfying the desired properties:
$\lim_{n \to \infty}\|f_{n} - f\|_{p} = 0$ and
$\lim(Tf_{n})(x) = (Tf^{0})(x) + (Tf^{1})(x) = (Tf)(x)$ almost everywhere.

We still need to remove the restrictions $\alpha > 0$ and $\beta < 1$. The
excluded cases, however, are simpler than the ones we have considered. If
$\alpha = 0$ and $\beta = 1$ then at least one of the pairs $(p_{0}, q_{0})$
and $(p_{1}, q_{1})$ must be $(\infty, 1)$ and there is nothing to prove. If
$\alpha > 0$ and $\beta = 1$ then the above proof is still valid if we define
$g_{z}$ to be $g$ for all $z \in \mathbb{C}$. Similarly, if $\alpha = 0$ and
$\beta < 1$, putting $f_{z} = f$ allows us to carry out the argument we gave.
This completes the proof of Theorem~\ref{thm:c5-1-3}.

The M. Riesz convexity theorem is a powerful tool in analysis and, in addition
to \eqref{eq:c5-1-1} and \eqref{eq:c5-1-2}, many important inequalities can be
derived from it. In some cases, however, one encounters a family of operators
$\{T_{t}\}$, $0 \leq t \leq 1$, mapping $L^{p_{t}}(M, d\mu_{t})$ boundedly into
$L^{q_{t}}(N, d\nu_{t})$, where both the operators and the measures vary in a
``smooth'' way with $t$. In the fourth section we shall derive an extension of
Theorem~\ref{thm:c5-1-3} that can be applied to such families of operators. In
other cases one encounters (not necessarily linear) operators defined on
$L^{p_{t}}$-spaces, $0 \leq t \leq 1$, that are not of types $(p_{0}, q_{0})$
and $(p_{1}, q_{1})$, yet they do map $L^{p_{t}}$ boundedly into $L^{q_{t}}$
when $0 < t < 1$. In the next section we shall show that certain weaker
conditions at the ``end points'' $(p_{i}, q_{i})$, $i = 0$, 1, are usually
sufficient in order to guarantee (strong) type $(p_{t}, q_{t})$ for
$0 < t < 1$.

\section{The Marcinkiewicz Interpolation Theorem}\index{Marcinkiewicz interpolation theorem@Marcinkiewicz interpolation theorem}

In the second chapter we introduced the maximal function of Hardy and
Littlewood, $m_{f}$, of a function $f \in L^{p}(E_{n})$, $1 \leq p \leq \infty$.
We showed (see Corollary \ref{thm:c2-3-11} of Chapter II) that the operator
$M$ mapping $f$ into $m_{f}$ is of type $(p, p)$ for $1 < p \leq \infty$. The
proof of this result depended on two ``end point'' results: One is the easily
established fact that $M$ is of type $(\infty, \infty)$ with norm 1; the other
result is Theorem~\ref{thm:c2-3-10} which describes the behavior of $M$
restricted to $L^{1}(E_{n})$. When $f$ belongs to $L^{1}(E_{n})$ we showed that
the distribution
function\footnote{If $(M, \mathcal{M}, \mu)$ is a measure space and $g$ is a
measurable function on $M$, its distribution function $\lambda$ is defined by
letting its value at a nonnegative real number $s$ be
$\lambda(s) = \mu\{x \in M ; |g(x)| > s\}$. In the third section of Chapter II
we introduced this concept when $M = E_{n}$ and $\mu$ is Lebesgue measure.},
$\lambda$, of $m_{f}$ satisfies the inequality
\begin{equation}\label{eq:c5-2-1}
\lambda(s)\leq\frac{c\|f\|_{1}}{s}
\end{equation}
for all $s > 0$, where the constant $c$ is independent of $f$. If there existed
an integrable function $\varphi$ on $(0, \infty)$ such that
$\lambda(s) \leq \varphi(s)\|f\|_{1}$, in view of the equality
$\|m_{f}\|_{1} = \int_{0}^{\infty}\lambda(s)\,ds$ (see \eqref{eq:c2-3-9} of
Chapter II), we could conclude that the operator $M$ is of type $(1, 1)$. But
\eqref{eq:c5-2-1} cannot be improved for the general function in $L^{1}(E_{n})$
(this can be easily seen by calculating $m_{f}$ when $f$ is the characteristic
function of the interval $(0, 1) \subset E_{1}$).

The argument that was used to show that $M$ is of type $(p, p)$,
$1 < p \leq \infty$, can be extended to a general theorem, the Marcinkiewicz
interpolation theorem, concerning operators mapping $L^{p}$ spaces into $L^{q}$
spaces. This result extends the M. Riesz convexity theorem in the sense
mentioned at the end of Section 1 (that weaker end-point conditions still
guarantee boundedness for intermediate values, and that operators that are
somewhat more general than linear are involved); on the other hand, it does not
apply to all the pairs $(p, q)$ considered in Theorem~\ref{thm:c5-1-3}, and we
do not obtain similar conclusions concerning the convexity of the norms.

In order to state the Marcinkiewicz theorem we need to introduce some new
concepts. Suppose $T$ is an operator defined on a linear space $D$ of
measurable functions on $(M, \mathcal{M}, \mu)$ having values that are
measurable functions defined on another measure space $(N, \mathcal{N}, \nu)$.
We say that $T$ is subadditive\index{subadditive (operator)@subadditive (operator)} if
\begin{equation}\label{eq:c5-2-2}
|[T(f_{1}+f_{2})](x)|\leq|[Tf_{1}](x)|+|[Tf_{2}](x)|
\end{equation}
for almost every $x$ in $N$, whenever $f_{1}$ and $f_{2}$ belong to $D$. If $D$
contains all finite linear combinations of characteristic functions of sets of
finite measure and contains all the truncations of its members, we say that $T$
is of weak type\index{weak type (p,q)@weak type $(p,q)$} $(p, q)$, $1 \leq p \leq \infty$ and $1 \leq q < \infty$, if
there exists a constant $k$, independent of $f \in L^{p}(M) \cap D$, such that
\begin{equation}\label{eq:c5-2-3}
\lambda(s)\leq\left(\frac{k\|f\|_{p}}{s}\right)^{q},
\end{equation}
where $\lambda$ is the distribution function of $Tf$. When $q = \infty$
inequality \eqref{eq:c5-2-3} is replaced by the condition
$\|Tf\|_{q} \leq k\|f\|_{p}$. The least $k$ for which these inequalities are
valid is called the weak\index{weak (p,q) norm@weak $(p,q)$ norm} $(p, q)$ norm of $T$.

It is easy to see that an operator that is of type $(p, q)$ must be of weak
type $(p, q)$. When $q = \infty$ this is obvious. Suppose, then, that
$q < \infty$, $f \in L^{p}(M) \cap D$, $s > 0$ and
$\lambda(s) = \nu\{x \in N ; |(Tf)(x)| > s\} = \nu(E_{s})$ is the distribution
function of $Tf$; then
\[
s^{q}\lambda(s)=s^{q}\int_{E_{s}}d\nu\leq\int_{E_{s}}|Tf|^{q}\,d\nu
\leq\|Tf\|_{q}^{q}\leq(k\|f\|_{p})^{q}
\]
and we obtain \eqref{eq:c5-2-3} by dividing by $s^{q}$.

The Marcinkiewicz interpolation theorem can now be formulated in the following
way:

\begin{theorem}\label{thm:c5-2-4}
Suppose $T$ is a subadditive operator of weak type $(p_{j}, q_{j})$, where
$1 \leq p_{j} \leq q_{j} \leq \infty$ for $j = 0$, 1 and $q_{0} \neq q_{1}$,
then $T$ is of type $(p_{t}, q_{t})$ provided
\[
\frac{1}{p_{t}}=\frac{1-t}{p_{0}}+\frac{t}{p_{1}}
\qquad\text{and}\qquad
\frac{1}{q_{t}}=\frac{1-t}{q_{0}}+\frac{t}{q_{1}},
\]
when $0 < t < 1$.
\end{theorem}

In the next section we shall introduce a class of function spaces (that
includes the $L^{p}$ spaces) which forms a natural setting for an extension of
Theorem~\ref{thm:c5-2-4}. Since we will give a proof of this extension in
Section 3 we do not prove the Marcinkiewicz theorem here. Instead, we devote
the rest of this section to an important application of this interpolation
theorem.

In the third chapter we introduced the Hilbert transform and mentioned some of
its properties (see \eqref{frp:c3-6-14} in Chapter III). In particular, we
stated that it is an operator of type $(p, p)$, $1 < p < \infty$. We shall show
that this result follows from the Marcinkiewicz theorem by first establishing
that it is of type $(2, 2)$ and then showing that it is of weak type $(1, 1)$.
Our previous treatment of this operator was very brief; consequently, we need
to develop some of its properties further.

Suppose $f \in L^{p}(-\infty, \infty)$, $1 \leq p < \infty$, and is
real-valued. Then, by Theorems~\ref{thm:c2-2-1} and~\ref{thm:c2-3-16} of the
second chapter, the Poisson integral of $f$
\[
u(x+iy)=u(x,y)=\frac{1}{\pi}\int_{-\infty}^{\infty}f(t)
\frac{y}{(x-t)^{2}+y^{2}}\,dt
\]
is a harmonic function in the upper half-plane $E_{2}^{+}$ such that
$u(z) = u(x + iy)$ tends to $f(x_{0})$ nontangentially for almost all
$x_{0} \in (-\infty, \infty)$ and
\begin{equation}\label{eq:c5-2-5}
\left(\int_{-\infty}^{\infty}|u(x+iy)|^{p}\,dx\right)^{1/p}\leq\|f\|_{p}
\end{equation}
for all $y > 0$.

It is not hard to show that there exists a (unique) harmonic function $v$,
defined on $E_{2}^{+}$, such that $F = u + iv$ is analytic in this region and
$\lim_{y \to \infty}v(x + iy) = 0$ for all $x \in (-\infty, \infty)$. This
function can be obtained by convolving $f$ with the conjugate Poisson kernel
\[
Q(x,y)=\frac{1}{\pi}\frac{x}{x^{2}+y^{2}},
\]
$-\infty < x < \infty$, $y > 0$ (see \eqref{frp:c3-6-13}, Chapter III). That is,
\[
\begin{aligned}&v(x,y)=v(x+iy)=\frac{1}{\pi}\int_{-\infty}^{\infty}f(t)
\frac{x-t}{(x-t)^{2}+y^{2}}\,dt\\
&\qquad=\frac{1}{\pi}\int_{-\infty}^{\infty}f(x-t)\frac{t}{t^{2}+y^{2}}\,dt.\end{aligned}
\]

Since
\[
P(x,y)+iQ(x,y)=\frac{1}{\pi}\left\{\frac{y}{x^{2}+y^{2}}+i\frac{x}{x^{2}+y^{2}}\right\}
=\frac{i}{\pi}\frac{1}{x+iy}=\frac{i}{\pi z}
\]
is analytic in $E_{2}^{+}$, it follows immediately that so is
$F(z) = u(z) + iv(z)$. The fact that $v(x + iy)$ tends to 0 as $y \to \infty$
is obvious.

\begin{lemma}\label{lem:c5-2-6}
The function $F(z)$, $z = x + iy \in E_{2}^{+}$, has a nontangential limit
$f(x_{0}) + i\tilde{f}(x_{0})$ at almost every point
$x_{0} \in (-\infty, \infty)$.
\end{lemma}

\begin{proof}
By considering the negative and positive parts of $f$ separately we can reduce
the problem to the case $f \leq 0$. Thus, $u \leq 0$ and the function
$G = \exp\{u + iv\} = e^{F}$ has absolute value $e^{u} \leq 1$. By
Theorem~\ref{thm:c2-2-5} of Chapter II, therefore, $G(z)$ has nontangential
limits a.e. on the boundary. These limits cannot be 0 on a set of positive
measure since $u$ is a Poisson integral. Thus, $v(z)$ must have nontangential
limits modulo $2\pi$ at almost all boundary points
$x_{0} \in (-\infty, \infty)$. That is, the limit points of $v(z)$, as $z$
tends nontangentially to $x_{0}$, have the form $a + 2k\pi$, where $k$ is an
integer. But it is an immediate consequence of the continuity of $v(z)$,
$z \in E_{2}^{+}$, that, if the values $v(z)$ have two distinct limit points,
$a$ and $b$, as $z$ approaches $x_{0}$ nontangentially, then all points between
$a$ and $b$ must also be limit points of this set of values. Hence, $v(z)$ can
have a nontangential limit modulo $2\pi$ only if it has a nontangential limit.
This proves the lemma.
\end{proof}

If, as in Lemma~\ref{lem:c5-2-6}, we let $\tilde{f}(x_{0})$ denote the
nontangential limit of $v(z)$ we obtain an almost everywhere defined function.
The mapping $f \to \tilde{f}$ (which we have shown to be defined for all
$f \in L^{p}(-\infty, \infty)$, $1 \leq p < \infty$) is called the Hilbert
transform\index{Hilbert transform@Hilbert transform}. In the third chapter (see \eqref{frp:c3-6-13} in Chapter III) we
indicated how one can show that the Hilbert transform is of type $(2, 2)$. We
first observed that, for each $y > 0$,
\[
\hat{Q}_{y}(t)=\int_{-\infty}^{\infty}e^{-2\pi itx}Q(x,y)\,dx
=(-i\operatorname{sgn}t)e^{-2\pi|yt|}.
\]
Thus, since $v(\cdot, y)$ is the convolution of $Q(\cdot, y)$ and $f$,
\[
\hat{v}_{y}(t)=\int_{-\infty}^{\infty}e^{-2\pi itx}v(x,y)\,dx
=(-i\operatorname{sgn}t)e^{-2\pi|yt|}\hat{f}(t)
\]
for a.e. $t$. By the Plancherel theorem and Fatou's lemma, therefore,
\begin{align*}
\|\tilde{f}\|_{2}^{2}&=\int_{-\infty}^{\infty}|\tilde{f}(x)|^{2}\,dx
\leq\liminf_{y \to 0}\int_{-\infty}^{\infty}|v(x+iy)|^{2}\,dx\\
&=\liminf_{y \to 0}\int_{-\infty}^{\infty}|\hat{v}_{y}(t)|^{2}\,dt
\leq\int_{-\infty}^{\infty}|\hat{f}(t)|^{2}\,dt
=\int_{-\infty}^{\infty}|f(t)|^{2}\,dt=\|f\|_{2}^{2}.
\end{align*}
This is the desired inequality showing that the Hilbert transform is of type
$(2, 2)$:
\begin{equation}\label{eq:c5-2-7}
\|\tilde{f}\|_{2}\leq\|f\|_{2}
\end{equation}
for all $f \in L^{2}(-\infty, \infty)$.\footnote{A modification of the argument
shows that inequality \eqref{eq:c5-2-7} can be reversed (thus the Hilbert
transform is a unitary operator on $L^{2}(-\infty, \infty)$). This question
will be considered again in a more general context in Chapter VI when we
introduce the Riesz transforms.}

\begin{lemma}\label{lem:c5-2-8}
The Hilbert transform is of weak-type $(1, 1)$.
\end{lemma}

\begin{proof}
Considering the positive and negative parts of $f \in L^{1}(-\infty, \infty)$
separately, we can reduce the problem to $f \geq 0$. As before, we put
\begin{equation}\label{eq:c5-2-9}
\begin{aligned}F(z)=F(x+iy)&=u(x,y)+iv(x,y)\\
&=\int_{-\infty}^{\infty}f(t)\{P(x-t,y)+iQ(x-t,y)\}\,dt\end{aligned}
\end{equation}
for $z \in E_{2}^{+}$. Then, for $s > 0$,
$w(x, y) = \log|1 + sF(z)|$ is harmonic in $E_{2}^{+}$ and bounded in each
proper sub-half-space $\{x + iy ; y \geq y_{0} > 0\}$. Thus, by
Lemma~\ref{lem:c2-2-7} of Chapter II,
\begin{align*}
w(x,y)&=\pi\log|1+sF(x+iy)|\\
&=\int_{-\infty}^{\infty}\frac{(y-\eta)}{(x-\xi)^{2}+(y-\eta)^{2}}
\log|1+sF(\xi+i\eta)|\,d\xi
\end{align*}
whenever $0 < \eta < y$. Letting $\eta \to 0$ and applying
Lemma~\ref{lem:c5-2-6} together with Fatou's lemma we see that the last
integral is greater than or equal to
\[
\int_{-\infty}^{\infty}\frac{y}{(x-\xi)^{2}+y^{2}}
\log\sqrt{(1+sf(\xi))^{2}+(s\tilde{f}(\xi))^{2}}\,d\xi.
\]
Consequently, multiplying by $y$ and using elementary properties of the
logarithm function we obtain
\begin{align*}
\frac{1}{\pi}\int_{-\infty}^{\infty}&\frac{y^{2}}{(x-\xi)^{2}+y^{2}}
\log\sqrt{(1+sf(\xi))^{2}+(s\tilde{f}(\xi))^{2}}\,d\xi\\
&\leq y\log|1+sF(x+iy)|\leq y\log(1+s|F(x+iy)|)\leq ys|F(x+iy)|.
\end{align*}
On the other hand, it follows immediately from the representation
\eqref{eq:c5-2-9} that $\pi yF(x + iy)$ tends to
$\|f\|_{1} = \int_{-\infty}^{\infty}f(\xi)\,d\xi$ as $y$ tends to $\infty$.
Thus, letting $y$ tend to $\infty$ in the first and last terms of the above
inequality we have
\begin{equation}\label{eq:c5-2-10}
\int_{-\infty}^{\infty}\log\sqrt{(1+sf(\xi))^{2}+(s\tilde{f}(\xi))^{2}}\,d\xi
\leq s\|f\|_{1}.
\end{equation}
If $E_{\tau} = \{\xi \in (-\infty, \infty) : |\tilde{f}(\xi)| > \tau\}$,
$\tau > 0$, and $m(E_{\tau})$ denotes the Lebesgue measure of $E_{\tau}$,
\eqref{eq:c5-2-10} implies
\[
\begin{aligned}&(\log s\tau)m(E_{\tau})\leq\int_{E_{\tau}}\log|s\tilde{f}(\xi)|\,d\xi\\
&\qquad\leq\int_{-\infty}^{\infty}\log\sqrt{(1+sf(\xi))^{2}+(s\tilde{f}(\xi))^{2}}\,d\xi\\
&\qquad\leq s\|f\|_{1}.\end{aligned}
\]
We now choose $s = e/\tau$ and obtain the desired weak type inequality
\[
m(E_{\tau})\leq e\|f\|_{1}/\tau.
\]
\end{proof}

Theorem~\ref{thm:c5-2-4} can now be applied to the Hilbert transform with
$p_{0} = 1 = q_{0}$ and $p_{1} = 2 = q_{1}$ and we obtain the fact that there
exists a constant $A_{p}$, independent of $f \in L^{p}(-\infty, \infty)$,
$1 < p \leq 2$, such that $\|\tilde{f}\|_{p} \leq A_{p}\|f\|_{p}$. It is not
hard to see that this inequality is also valid when $2 \leq p < \infty$. First
we observe that the Hilbert transform, being a limit of convolution operators,
commutes with translations. The inequality
\begin{equation}\label{eq:c5-2-11}
\|\tilde{f}\|_{p}\leq A_{p}\|f\|_{p},
\end{equation}
$f \in L^{p}(-\infty, \infty)$, $1 < p < \infty$, is then a consequence of the
just established fact that it is valid when $1 < p \leq 2$ and
Theorem~\ref{thm:c1-3-20} of Chapter I.

\section{\texorpdfstring{$L(p,q)$}{L(p,q)} Spaces}\index{L(p,q) spaces@$L(p,q)$ spaces}

An operator of weak type $(r, p)$ is one that maps an $L^{r}$ space into a
class of functions $f$ having distribution functions $\lambda$ satisfying
\begin{equation}\label{eq:c5-3-1}
\sup_{s>0}\Phi(s)=\sup_{s>0}\{s[\lambda(s)]^{1/p}\}=A<\infty.
\end{equation}
This inequality can be restated by saying that the $L^{\infty}$ norm of
$\Phi$ is finite. We have already observed that
\begin{equation}\label{eq:c5-3-2}
\|f\|_{p}=\left(\int_{0}^{\infty}s^{p}\,d\{-\lambda(s)\}\right)^{1/p}
=\left(p\int_{0}^{\infty}s^{p-1}\lambda(s)\,ds\right)^{1/p}
\end{equation}
(see \eqref{eq:c2-3-6} in Chapter II). Thus, $f$ belongs to $L^{p}$ if and only
if $\Phi$ belongs to $L^{p}(0, \infty)$ with respect to the measure
$d\nu(s) = [1/\lambda(s)]\,d\{-\lambda(s)\}$. In fact, we have
\[
\|f\|_{p}=\left(\int_{0}^{\infty}[\Phi(s)]^{p}\,d\nu(s)\right)^{1/p}.
\]

Suppose that, instead of the $L^{\infty}$-norm or the $L^{p}$-norm of $\Phi$,
we considered an $L^{q}$-norm, $1 \leq q < \infty$; that is, we imposed the
condition
\begin{equation}\label{eq:c5-3-3}
\left(\int_{0}^{\infty}[\Phi(s)]^{q}\,d\nu(s)\right)^{1/q}<\infty.
\end{equation}
We could then ask whether operators mapping $L^{r}$ spaces into classes of
functions satisfying \eqref{eq:c5-3-3} can be interpolated. In this section we
shall study function spaces defined by conditions of this type and shall show
that they provide a natural setting for the Marcinkiewicz theorem.

The expression occurring in \eqref{eq:c5-3-3} does not define a norm.
Nevertheless, there are inequalities equivalent to \eqref{eq:c5-3-3} that, in
most cases, involve the finiteness of a\index{non-increasing rearrangement@non-increasing rearrangement!of a function@of a function} norm. In order to obtain these
conditions we introduce the nonincreasing rearrangement $f^{*}$ of a measurable
function $f$ defined on the measure space $(M, \mathcal{M}, \mu)$. If $\lambda$
denotes the distribution function of $f$, then $f^{*}$ is defined for $t > 0$
by letting
\[
f^{*}(t)=\inf\{s : \lambda(s)\leq t\}.
\]
It is clear that the function $\lambda$ is
nonincreasing\footnote{Since $\{x \in M : |f(x)| > s_{1}\} \subset \{x \in M :
|f(x)| > s_{2}\}$ when $s_{1} > s_{2}$ and the measure $\mu$ is monotone.}; if
it were a continuous and strictly decreasing function onto the positive reals
it would have an inverse function having these same properties. It follows
immediately from the above definition that, in this case, $f^{*}$ is this
inverse function. The following series of lemmas are devoted to establishing
some of the basic properties of the functions $\lambda$ and $f^{*}$. In the
sequel we shall assume that $\lim_{s \to \infty}\lambda(s) = 0$ and, hence,
$f^{*}(t) < \infty$ for all $t > 0$.

\begin{lemma}\label{lem:c5-3-4}
\textup{(i)} $\lambda$ and $f^{*}$ are nonincreasing and continuous on the
right,

\textup{(ii)} $\lambda(f^{*}(t)) \leq t$ for all $t > 0$;

\textup{(iii)} $f$ and $f^{*}$ have the same distribution function.\index{f*@$f^{*}$}
\end{lemma}

\begin{proof}
The fact that $\lambda$ is nonincreasing was observed above. It is also
immediate from its definition that $f^{*}$ must also be nonincreasing.
Continuity on the right for $\lambda$ follows from the continuity from below
of the measure $\mu$ and the fact that
\[
\lim_{s>s_{0},\, s \to s_{0}}\{x \in M : |f(x)| > s\}=\{x \in M : |f(x)| > s_{0}\}.
\]
Inequality (ii) follows immediately from this continuity on the right.

In order to show that $f^{*}$ is continuous on the right at $t_{0}$ we first
observe that this is immediate if $f^{*}(t_{0}) = 0$ (since, in this case,
$f^{*}(t) = 0$ for $t > t_{0}$ because $f^{*}$ is nonnegative and
nonincreasing). If $f^{*}(t_{0})$ is positive let $\alpha$ be such that
$f^{*}(t_{0}) > \alpha > 0$ and $\{\varepsilon_{n}\}$ a sequence of positive
real numbers decreasing to 0. From the definition of $f^{*}$ we must have
$\lambda(f^{*}(t_{0}) - \alpha) > t_{0}$. Thus, there exists $n_{0}$ such that
$\lambda(f^{*}(t_{0}) - \alpha) > t_{0} + \varepsilon_{n}$ for $n \geq n_{0}$.
But this implies $f^{*}(t_{0}) - \alpha < f^{*}(t_{0} + \varepsilon_{n})$ for
$n \geq n_{0}$; otherwise, for some $n \geq n_{0}$,
$f^{*}(t_{0}) - \alpha \geq f^{*}(t_{0} + \varepsilon_{n})$ and, by the
nonincreasing property of $\lambda$ and (ii), we would have the contradiction
$\lambda(f^{*}(t_{0}) - \alpha) \leq \lambda(f^{*}(t_{0} + \varepsilon_{n})) \leq t_{0} + \varepsilon_{n}$.
Thus, since $f^{*}$ is nonincreasing
$f^{*}(t_{0}) - \alpha < f^{*}(t_{0} + \varepsilon_{n}) \leq f^{*}(t_{0})$ for
all $n \geq n_{0}$. This shows that $f^{*}$ is continuous on the right.

In order to show property (iii) we first observe that it follows from the
definition of $f^{*}$ that $f^{*}(t) > s$ if and only if $t < \lambda(s)$.
Thus, $E_{s}^{*} = \{t > 0 : f^{*}(t) > s\}$ is precisely the interval
$(0, \lambda(s))$. Property (iii) now follows from the fact that the value of
the distribution function of $f^{*}$ at $s$ is the Lebesgue measure of the set
$E_{s}^{*}$.
\end{proof}

\begin{lemma}\label{lem:c5-3-5}
Suppose $\{f_{m}\}$ is a sequence of measurable functions such that, for all
$x \in M$, $|f_{m}(x)| \leq |f_{m+1}(x)|$, $m = 1, 2, 3, \ldots$. If $f$ is a
measurable function satisfying $|f(x)| = \lim_{m \to \infty}|f_{m}(x)|$ for
$x \in M$ then

\textup{(i)} for each $s > 0$, $\lambda_{m}(s)$ increases monotonically to
$\lambda(s)$ as $m$ tends to $\infty$, where $\lambda_{m}$ and $\lambda$ are
the distribution functions of $f_{m}$ and $f$;

\textup{(ii)} for each $t > 0$, $f_{m}^{*}(t)$ increases monotonically to
$f^{*}(t)$.
\end{lemma}

\begin{proof}
It is clear that
\[
E_{s}^{(m)}=\{x \in M : |f_{m}(x)| > s\}\subset E_{s}=\{x \in M : |f(x)| > s\}
\]
and $\bigcup_{m=1}^{\infty}E_{s}^{(m)} = E_{s}$. Thus, since the measure $\mu$
is monotone and continuous from below we must have
$\lambda_{m}(s) = \mu(E_{s}^{(m)}) \leq \mu(E_{s}) = \lambda(s)$ and
$\lim_{m \to \infty}\lambda_{m}(s) = \lambda(s)$. This proves (i). It now
follows from the definition of the nonincreasing rearrangement that
$f_{m}^{*}(t) \leq f_{m+1}^{*}(t) \leq f^{*}(t)$ for $m = 1, 2, 3, \ldots$.
Let $l = \lim_{m \to \infty}f_{m}^{*}(t)$. Since $f_{m}^{*}(t) \leq l$ we must
have $\lambda_{m}(l) \leq \lambda_{m}(f_{m}^{*}(t)) \leq t$ (the last
inequality being a consequence of part (ii) of Lemma~\ref{lem:c5-3-4}). Thus,
$\lambda(l) = \lim_{m \to \infty}\lambda_{m}(l) \leq t$ and this implies
$f^{*}(t) \leq l$. But, from the inequality $f_{m}^{*}(t) \leq f^{*}(t)$ we
obtain $l \leq f^{*}(t)$. It therefore follows that $l = f^{*}(t)$ and the
lemma is proved.
\end{proof}

This lemma is particularly useful in proving other properties of the functions
$\lambda$ and $f^{*}$ since it allows us to reduce these proofs to the case
when $f$ is a simple function. The distribution function and the nonincreasing
rearrangement of such a function is particularly simple to describe. Suppose
$f = \sum_{j=1}^{n}c_{j}\chi_{E_{j}}$, where $E_{j} \in \mathcal{M}$,
$\mu(E_{j}) > 0$ and $E_{j} \cap E_{k} = \emptyset$ if $j \neq k$. By
relabelling, if necessary, we can assume $c_{1} > c_{2} > \cdots > c_{n} >
c_{n+1} = 0$. Let $d_{j} = \mu(E_{1}) + \cdots + \mu(E_{j})$,
$1 \leq j \leq n$, and define $d_{0}$ to be 0. Then, the distribution function
$\lambda$ of $f$ has the form
\begin{equation}\label{eq:c5-3-6}
\lambda(s)=\begin{cases}
d_{j} & \quad\text{if } c_{j+1} \leq s < c_{j},\ 1 \leq j \leq n,\\
0 & \quad\text{if } c_{1} \leq s.\\
\end{cases}
\end{equation}
It follows that
\begin{equation}\label{eq:c5-3-7}
f^{*}(t)=\begin{cases}
c_{j} & \quad\text{if } d_{j-1} \leq t < d_{j},\ 1 \leq j \leq n,\\
0 & \quad\text{if } d_{n} \leq t.\\
\end{cases}
\end{equation}
By inspection, therefore, we see that for $p > 0$
\[
\sup_{s>0}s[\lambda(s)]^{1/p}=\sup_{1 \leq j \leq n}d_{j}^{1/p}c_{j}
=\sup_{t>0}t^{1/p}f^{*}(t).
\]

From this we can easily show that the supremum in \eqref{eq:c5-3-1} can be
expressed in terms of the nonincreasing rearrangement of the function in
question:

\begin{lemma}\label{lem:c5-3-8}
If $f$ is measurable, $\lambda$ is its distribution function and $p > 0$ then
\eqref{eq:c5-3-1} holds if and only if
\[
\sup_{t>0}t^{1/p}f^{*}(t)=A<\infty.
\]
\end{lemma}

\begin{proof}
We have just seen that Lemma~\ref{lem:c5-3-8} is true when $f$ is simple. If
$f$ is only measurable we can find a sequence $\{f_{m}\}$ of nonnegative simple
functions that increase monotonically to $|f|$ at all the points of $M$; the
lemma then follows immediately from Lemma~\ref{lem:c5-3-5}.
\end{proof}

In view of inequality \eqref{eq:c5-3-3} and the comments made before it, it is
natural to consider the space $L(p, q)$ of all those measurable functions $f$
satisfying
\[
\|f\|_{pq}^{*}=\left(\frac{q}{p}\int_{0}^{\infty}[t^{1/p}f^{*}(t)]^{q}\frac{dt}{t}\right)^{1/q}<\infty,
\]
when $1 \leq p < \infty$, $1 \leq q < \infty$, and
\[
\|f\|_{pq}^{*}=\sup_{t>0}t^{1/p}f^{*}(t)<\infty
\]
when $1 \leq p \leq \infty$ and $q = \infty$.\footnote{These definitions make
sense also when $0 < p < 1$ and $0 < q < 1$; however, since we shall be
interested in converting $L(p, q)$ into a Banach space, we do not consider
these cases. Moreover, in the above definition we adopt the usual convention
that $1/\infty = 0$. Observe that we have defined $L(\infty, q)$ only when
$q = \infty$.}

When $f$ is simple we see immediately from \eqref{eq:c5-3-7} that
\begin{equation}\label{eq:c5-3-9}
\begin{aligned}\|f\|_{p}&=\left(\int_{M}|f|^{p}\,d\mu\right)^{1/p}\\
&=\left(\int_{0}^{\infty}[f^{*}(t)]^{p}\,dt\right)^{1/p}=\|f^{*}\|_{p},
\qquad 1 \leq p < \infty,\end{aligned}
\end{equation}
and
\begin{equation}\label{eq:c5-3-9p}
\tag{3.9$'$}
\|f\|_{\infty}=\|f^{*}\|_{\infty}.
\end{equation}
From Lemma~\ref{lem:c5-3-5} we obtain the fact that these equalities hold for
all $f \in L^{p}(M)$, $1 \leq p \leq \infty$. Since
$\|f\|_{pp}^{*} = (\int_{0}^{\infty}[t^{1/p}f^{*}(t)]^{p}\,(dt/t))^{1/p} = (\int_{0}^{\infty}[f^{*}(t)]^{p}\,dt)^{1/p} = \|f^{*}\|_{p}$
it follows that $L(p, p) = L^{p}(M)$ and $\|\ \|_{pp}^{*}$ is a norm. In
general, however, $\|\ \|_{pq}^{*}$ is not a norm since the Minkowski
inequality may fail. Despite this, we will show that $\|\ \|_{pq}^{*}$ can be
used in order to introduce a topology on $L(p, q)$ which, in most cases, can be
shown to be that of a Banach space.

There are several reasons for introducing the constant $q/p$ in the definition
of $\|\ \|_{pq}^{*}$. One of the consequences of this definition is the fact
that $\|\chi_{E}\|_{pq}^{*}$ is independent of $q$ whenever $E$ is a measurable
subset of $M$ of finite measure and $\chi_{E}$ its characteristic function. In
fact, we have
\begin{equation}\label{eq:c5-3-10}
\|\chi_{E}\|_{pq}^{*}=\{\mu(E)\}^{1/p}
\end{equation}
for $1 \leq q \leq \infty$. This equality is a reflection of the fact that each
$L(p, q)$ space, for $p$ fixed, is a ``type of $L^{p}$ space.'' The following
theorem gives us a precise formulation of this assertion. It also shows that,
in some sense, $L(p, 1)$ is the ``smallest'' normed linear space for which
\eqref{eq:c5-3-10} holds, while $L(p, \infty)$ is the ``largest'' such space.

\begin{theorem}\label{thm:c5-3-11}
If $f \in L(p, q_{1})$ and $q_{1} \leq q_{2}$ then
\[
\textup{(i)}\qquad \|f\|_{pq_{2}}^{*}\leq\|f\|_{pq_{1}}^{*};
\]
consequently,
\[
\textup{(ii)}\qquad L(p,q_{1})\subset L(p,q_{2}).%
\footnote{In order to obtain this inclusion relation it suffices to show that
there exists a constant $B$, independent of $f \in L(p, q_{1})$, such that
$\|f\|_{pq_{2}}^{*} \leq B\|f\|_{pq_{1}}^{*}$. Short proofs establishing this
inequality can be found. For example:
$(\int_{0}^{\infty}[t^{1/p}f^{*}(t)]^{q_{2}}\,dt/t)^{q_{1}/q_{2}}
\leq (\sum_{k=-\infty}^{\infty}[f^{*}(2^{k-1})]^{q_{2}}\int_{2^{k-1}}^{2^{k}}t^{(q_{2}/p)-1}\,dt)^{q_{1}/q_{2}}
\leq B'\sum_{k=-\infty}^{\infty}[f^{*}(2^{k-1})]^{q_{1}}2^{kq_{1}/p}
\leq B''\sum_{k=-\infty}^{\infty}\int_{2^{k-2}}^{2^{k-1}}[t^{1/p}f^{*}(t)]^{q_{1}}\,dt/t
\leq B''\{\|f\|_{pq_{1}}^{*}\}^{q_{1}}$.}
\]

Suppose $\|\ \|$ is an order preserving norm (we mean by this that
$\|g\| \leq \|f\|$ whenever $|g(x)| \leq |f(x)|$ a.e.) defined on the simple
functions on $M$; then

\textup{(iii)} $\|\chi_{E}\| \leq \{\mu(E)\}^{1/p}$ for all $E \in \mathcal{M}$
implies $\|f\| \leq \|f\|_{p1}^{*}$ for all simple functions $f$;

\textup{(iv)} $\{\mu(E)\}^{1/p} \leq \|\chi_{E}\|$ for all $E \in \mathcal{M}$
implies $\|f\|_{p\infty}^{*} \leq \|f\|$ for all simple functions $f$.
\end{theorem}

Inequality (i), however, is more subtle and requires a somewhat more difficult
proof.

Equality \eqref{eq:c5-3-10} shows that (i) is best possible.

\begin{proof}
Inequality (i) can be shown in a particularly simple way when $q_{2} = \infty$.
Since $f^{*}$ is nonincreasing,
\begin{align*}
t^{1/p}f^{*}(t)&=f^{*}(t)\left\{(q_{1}/p)\int_{0}^{t}u^{(q_{1}/p)-1}\,du\right\}^{1/q_{1}}\\
&\leq\left\{(q_{1}/p)\int_{0}^{t}[u^{1/p}f^{*}(u)]^{q_{1}}\frac{du}{u}\right\}^{1/q_{1}}
\leq\|f\|_{pq_{1}}^{*}.
\end{align*}
By taking the supremum of $t^{1/p}f^{*}(t)$, over all $t > 0$, we obtain the
desired inequality.

When $q_{2} < \infty$, it follows from Lemma~\ref{lem:c5-3-5} that it suffices
to establish the lemma for $f$ simple. We assume, therefore, that $f^{*}$ has
the form \eqref{eq:c5-3-7}, in which case
\[
\|f\|_{pq}^{*}=\left\{\sum_{j=1}^{n}c_{j}^{q}(d_{j}^{(q/p)}-d_{j-1}^{(q/p)})\right\}^{1/q}.
\]
Setting $a_{j} = c_{j}^{q_{2}}$, $b_{j} = d_{j}^{q_{2}/p}$ and
$\theta = q_{1}/q_{2}$, therefore, our inequality is equivalent to
\[
\textup{(a)}\qquad
\sum_{j=1}^{n}a_{j}(b_{j}-b_{j-1})\leq
\left\{\sum_{j=1}^{n}a_{j}^{\theta}(b_{j}^{\theta}-b_{j-1}^{\theta})\right\}^{1/\theta},
\]
where $a_{1} > a_{2} > \cdots > a_{n} > 0$, $0 = b_{0} < b_{1} < \cdots < b_{n}$
and $0 < \theta \leq 1$. We shall prove (a) by induction. When $n = 1$ this
inequality reduces to $a_{1}b_{1} \leq \{a_{1}^{\theta}b_{1}^{\theta}\}^{1/\theta}$,
which is obviously true.

Assume that (a) is true for $n = N$. We let
\[
\varphi(x)=\left\{\sum_{j=1}^{N}a_{j}^{\theta}(b_{j}^{\theta}-b_{j-1}^{\theta})
+x^{\theta}(b_{N+1}^{\theta}-b_{N}^{\theta})\right\}^{1/\theta}
=(A+x^{\theta}B)^{1/\theta}
\]
and
\[
l(x)=\sum_{j=1}^{N}a_{j}(b_{j}-b_{j-1})+x(b_{N+1}-b_{N})=A_{1}+xB_{1}
\qquad\text{for } 0 \leq x \leq a_{N}.
\]
We want to show that $\varphi(a_{N+1}) \geq l(a_{N+1})$ whenever
$0 < a_{N+1} < a_{N}$. By the induction assumption $\varphi(0) \geq l(0)$ and
$\varphi(a_{N}) \geq l(a_{N})$ (this last inequality reduces to (a) with $b_{N}$
replaced by $b_{N+1}$). On the other hand, since the derivative
$\varphi'(x) = B(Ax^{-\theta} + B)^{(1/\theta)-1}$ of $\varphi$ is decreasing on
the positive real axis, the function $\varphi$ must be concave. Therefore, the
fact that it dominates the linear function $l$ at the end points 0 and $a_{N}$
implies that $\varphi(x) \geq l(x)$ for $0 \leq x \leq a_{N}$. This proves part
(i). Part (ii) is then an immediate consequence.

Since $\|\ \|$ is order preserving it suffices to show (iii) when
$f = \sum_{j=1}^{n}c_{j}\chi_{E_{j}}$ is a nonnegative simple function. We can
assume that $c_{1} > c_{2} > \cdots > c_{n} > c_{n+1} = 0$ and the sets
$E_{1}, \ldots, E_{n}$ are pairwise disjoint. If $k = 1, 2, \ldots, n$ define
$f_{k} = b_{k}\chi_{F_{k}}$, where $F_{k} = \bigcup_{j=1}^{k}E_{j}$ and
$b_{k} = c_{k} - c_{k+1}$. Then, clearly,
\[
\textup{(b)}\qquad
f^{*}(t)=\sum_{k=1}^{n}f_{k}^{*}(t)\qquad\text{for } t > 0,
\]
and
\begin{align*}
\|f\|&\leq\sum_{k=1}^{n}\|f_{k}\|=\sum_{k=1}^{n}b_{k}\|\chi_{F_{k}}\|\\
&\leq\sum_{k=1}^{n}b_{k}\{\mu(F_{k})\}^{1/p}
=\sum_{k=1}^{n}\frac{1}{p}\int_{0}^{\infty}t^{1/p}f_{k}^{*}(t)\frac{dt}{t}\\
&=\frac{1}{p}\int_{0}^{\infty}t^{1/p}\sum_{k=1}^{n}f_{k}^{*}(t)\frac{dt}{t}
=\frac{1}{p}\int_{0}^{\infty}t^{1/p}f^{*}(t)\frac{dt}{t}=\|f\|_{p1}^{*}.
\end{align*}
This proves (iii).

In order to show (iv) we observe that, in terms of the notation developed
before equality \eqref{eq:c5-3-6}, we have
\[
\|f\|_{p\infty}^{*}=\sup_{1 \leq j \leq n}d_{j}^{1/p}c_{j}
\]
(this fact follows from the definition of $\|\ \|_{p\infty}^{*}$ and the
equalities displayed immediately after \eqref{eq:c5-3-7}). Suppose that the
supremum on the right is assumed at $j = k$; thus,
\[
\|f\|_{p\infty}^{*}=d_{k}^{1/p}c_{k}.
\]
Let $g = c_{k}\chi_{F_{k}}$. Then $0 \leq g \leq f$ and, thus,
\[
\|f\|_{p\infty}^{*}=d_{k}^{1/p}c_{k}=c_{k}\{\mu(F_{k})\}^{1/p}
\leq c_{k}\|\chi_{F_{k}}\|=\|g\|\leq\|f\|
\]
and the theorem is proved.
\end{proof}

It is useful to associate with each space $L(p, q)$ the point $(1/p, 1/q)$ of
the square $Q = \{(x_{1}, x_{2}) \in E_{2} : 0 \leq x_{1}, x_{2} \leq 1\}$. As
a consequence of Theorem~\ref{thm:c5-3-11}, part (ii), we see that if
$x_{1} = 1/p$ is fixed and $x_{2} = 1/q$ is allowed to vary from 1 to 0 we
obtain the points associated with an increasing family of function spaces. When
we reach the diagonal $\{(x_{1}, x_{2}) \in Q : x_{1} = x_{2}\}$ of the square
$Q$, the associated space is $L^{p}$ (see \eqref{eq:c5-3-9}). In view of the
definition of a weak type operator that was made in Section 2, it is natural to
refer to $L(p, \infty)$, the largest of the spaces encountered in this manner,
as weak\index{weak Lp@weak $L^{p}$} $L^{p}$. The justification for this terminology is that the inequality
\eqref{eq:c5-2-3} that was used in the definition of an operator of weak type
$(r, p)$ is, by Lemma~\ref{lem:c5-3-8}, equivalent to
\[
\|Tf\|_{p\infty}^{*}\leq k\|f\|_{r}.
\]
Since $\|f\|_{r} = \|f\|_{rr}^{*}$ and $r \geq 1$, it follows from inequality
(i) of Theorem~\ref{thm:c5-3-11} that $\|f\|_{r} \leq \|f\|_{r1}^{*}$. Thus, an
operator $T$ of weak type $(r, p)$ satisfies the weaker inequality
\begin{equation}\label{eq:c5-3-12}
\|Tf\|_{p\infty}^{*}\leq k\|f\|_{r1}^{*}
\end{equation}
for all $f$ belonging to the domain of $T$ and to $L(r, 1)$. The interpolation
theorem that will be proved in this section asserts, in part, that end point
conditions like \eqref{eq:c5-3-12}, even though weaker than weak type
conditions, are sufficient in order to guarantee the conclusion of the
Marcinkiewicz theorem (Theorem~\ref{thm:c5-2-4}). Closely connected with part
(iii) of Theorem~\ref{thm:c5-3-11} is the fact that the hypotheses of the
Marcinkiewicz theorem can be reduced still further. In fact, we have the
following result, showing that the conclusion of Theorem~\ref{thm:c5-2-4}
depends only on knowing that the weak type inequalities at the end points hold
for characteristic functions of sets of finite measure.

\begin{theorem}\label{thm:c5-3-13}
Suppose $T$ is a linear operator which maps the finite linear combinations of
characteristic functions $\chi_{E}$ of sets $E \subset M$ of finite measure
into a vector space $B$ that is endowed with an order preserving norm
$\|\cdot\|$. If
\[
\|T\chi_{E}\|\leq C\|\chi_{E}\|_{r1}^{*}=C\{\mu(E)\}^{1/r},
\]
where $C$ is independent of $E$, then there exists a constant $A$ such that
\[
\|Tf\|\leq A\|f\|_{r1}^{*}
\]
for all $f$ in the domain of $T$.
\end{theorem}

\begin{proof}
If $f \geq 0$ belongs to the domain of $T$ we can represent it as a sum
$f = \sum_{k=1}^{n}f_{k}$, where the functions $f_{k}$ are the multiples of
characteristic functions defined in the proof of Theorem~\ref{thm:c5-3-11} so
that equality (b) is satisfied. Then
\[
\|Tf\|=\left\|\sum_{k=1}^{n}Tf_{k}\right\|\leq\sum_{k=1}^{n}\|Tf_{k}\|
\leq\sum_{k=1}^{n}Cb_{k}\|\chi_{F_{k}}\|_{r1}^{*}.
\]
But, as was shown in the proof of Theorem~\ref{thm:c5-3-11}, part (iii), the
last sum equals $C\|f\|_{r1}^{*}$. If $f = f_{1} + if_{2}$ is complex valued we
obtain the theorem with $A = 4C$ by applying the result that has been just
established to the positive and negative parts of $f_{1}$ and
$f_{2}$.\footnote{Part (iii) of Theorem~\ref{thm:c5-3-11} is the special case
of Theorem~\ref{thm:c5-3-13} obtained when $T$ is the identity operator. It is
clear that the proof of Theorem~\ref{thm:c5-3-11} is valid for sublinear
operators $T$; that is, operators that are subadditive and positive homogeneous
($|Taf| = |a||Tf|$ for all scalars $a$).}
\end{proof}

The interpolation theorem we have been discussing is a consequence of the
following classical estimate, known as Hardy's\index{Hardy inequality@Hardy inequality} inequality.

\begin{lemma}\label{lem:c5-3-14}
If $q \geq 1$, $r > 0$ and $g$ is a nonnegative function defined on
$(0, \infty)$ then
\[
\textup{(i)}\qquad
\left(\int_{0}^{\infty}\left[\int_{0}^{t}g(u)\,du\right]^{q}t^{-r-1}\,dt\right)^{1/q}
\leq\frac{q}{r}\left(\int_{0}^{\infty}[ug(u)]^{q}u^{-r-1}\,du\right)^{1/q};
\]
\[
\textup{(ii)}\qquad
\left(\int_{0}^{\infty}\left[\int_{t}^{\infty}g(u)\,du\right]^{q}t^{r-1}\,dt\right)^{1/q}
\leq\frac{q}{r}\left(\int_{0}^{\infty}[ug(u)]^{q}u^{r-1}\,du\right)^{1/q}.
\]
\end{lemma}

\begin{proof}
We first show that (i) implies (ii). If we apply (i) to the function
$g_{1}(u) = u^{-2}g(u^{-1})$ we obtain that
\[
\begin{aligned}\int_{0}^{\infty}\left[\int_{0}^{t}g_{1}(u)\,du\right]^{q}t^{-r-1}\,dt
&=\int_{0}^{\infty}\left[\int_{1/t}^{\infty}g(v)\,dv\right]^{q}t^{-r-1}\,dt\\
&=\int_{0}^{\infty}\left[\int_{s}^{\infty}g(v)\,dv\right]^{q}s^{r-1}\,ds\end{aligned}
\]
is less than or equal to $(q/r)^{q}$ times
\[
\int_{0}^{\infty}[ug_{1}(u)]^{q}u^{-r-1}\,du
=\int_{0}^{\infty}[vg(v)]^{q}v^{r-1}\,dv.
\]
Thus, it suffices to establish inequality (i).

If we apply Jensen's inequality, with $\varphi(x) = |x|^{q}$ and
$d\mu(u) = u^{(r/q)-1}\,du$ (see Section 4, Example (2), in the second chapter)
we obtain
\begin{align*}
\left[\int_{0}^{t}g(u)\,du\right]^{q}
&=\left[\int_{0}^{t}g(u)u^{1-(r/q)}u^{(r/q)-1}\,du\right]^{q}\\
&\leq\left(\frac{q}{r}\right)^{q-1}t^{r(1-1/q)}
\int_{0}^{t}[g(u)]^{q}u^{q-r-1+r/q}\,du.
\end{align*}
Thus,
\begin{align*}
&\int_{0}^{\infty}\left[\int_{0}^{t}g(u)\,du\right]^{q}t^{-r-1}\,dt\\
&\quad\leq\left(\frac{q}{r}\right)^{q-1}\int_{0}^{\infty}t^{-1-r/q}
\left[\int_{0}^{t}[g(u)]^{q}u^{q-r-1+r/q}\,du\right]dt\\
&=\left(\frac{q}{r}\right)^{q-1}\int_{0}^{\infty}[g(u)u]^{q}u^{-r-1+r/q}
\left(\int_{u}^{\infty}t^{-1-r/q}\,dt\right)du\\
&=\left(\frac{q}{r}\right)^{q}\int_{0}^{\infty}[g(u)u]^{q}u^{-r-1}\,du
\end{align*}
and Hardy's inequality is proved.
\end{proof}

We can now state and prove the principal result of this section. We say that a
subadditive operator $T$ is of restricted weak\index{restricted weak type (p,q)@restricted weak type (p,q)} type $(r, p)$ if its domain $D$
contains all the truncations of its members as well as all finite linear
combinations of characteristic functions of sets of finite measure and satisfies
inequality \eqref{eq:c5-3-12} whenever $f \in D \cap L(r, 1)$. In view of
Theorem~\ref{thm:c5-3-13} such an operator may be considered to be one that
satisfies the weak type $(r, p)$ conditions when it is restricted to
characteristic functions of sets of finite measure.

\begin{theorem}\label{thm:c5-3-15}
Suppose $T$ is a subadditive operator of restricted weak types $(r_{j}, p_{j})$,
$j = 0$, 1, with $r_{0} < r_{1}$ and $p_{0} \neq p_{1}$, then there exists a
constant $B = B_{\theta}$ such that
\[
\|Tf\|_{pq}^{*}\leq B\|f\|_{rq}^{*}
\]
for all $f$ belonging to the domain of $T$ and to $L(r, q)$, where
$1 \leq q \leq \infty$,
\[
\frac{1}{p}=\frac{(1-\theta)}{p_{0}}+\frac{\theta}{p_{1}},\qquad
\frac{1}{r}=\frac{1-\theta}{r_{0}}+\frac{\theta}{r_{1}}
\qquad\text{and}\qquad 0<\theta<1.
\]
\end{theorem}

\begin{proof}
For $f \in L(r, q) \cap D$ we define
\[
f^{t}(x)=\begin{cases}
f(x) & \quad\text{if } |f(x)| > f^{*}(t^{\gamma}),\\
0 & \quad\text{if } |f(x)| \leq f^{*}(t^{\gamma}),\\
\end{cases}
\]
and $f_{t}(x) = f(x) - f^{t}(x)$, where
\[
\gamma=\frac{(1/p_{0})-(1/p)}{(1/r_{0})-(1/r)}
=\frac{(1/p)-(1/p_{1})}{(1/r)-(1/r_{1})}.
\]
We then have the two inequalities
\[
\textup{(a)}\qquad
(f^{t})^{*}(s)\leq\begin{cases}
f^{*}(s) & \quad\text{if } 0 < s < t^{\gamma},\\
0 & \quad\text{if } t^{\gamma} \leq s,\\
\end{cases}
\qquad\text{and}
\]
\[
f_{t}^{*}(s)\leq\begin{cases}
f^{*}(t^{\gamma}) & \quad\text{if } 0 < s < t^{\gamma},\\
f^{*}(s) & \quad\text{if } t^{\gamma} \leq s.\\
\end{cases}
\]
When $f$ is simple (a) can be derived readily from formulas analogous to
\eqref{eq:c5-3-7}; the general case then follows easily from this one. Since the
operator $T$ is subadditive
$|[Tf](y)| = |[T(f^{t} + f_{t})](y)| \leq |[Tf^{t}](y)| + |[Tf_{t}](y)|$ for
almost every $y \in N$. Thus, if $s > 0$,
\begin{align*}
&\{y \in N:{}|[Tf](y)|>(Tf^{t})^{*}(s)+(Tf_{t})^{*}(s)\}\\
&\quad\subset\{y \in N : |[Tf^{t}](y)|>(Tf^{t})^{*}(s)\}\\
&\quad\cup\{y \in N : |[Tf_{t}](y)|>(Tf_{t})^{*}(s)\}.
\end{align*}
If $\lambda$, $\lambda^{t}$, and $\lambda_{t}$ are the distribution functions of
$Tf$, $Tf^{t}$, and $Tf_{t}$, respectively, it follows from this inclusion
relation and Lemma~\ref{lem:c5-3-4}, part (ii),
\[
\lambda((Tf^{t})^{*}(s)+(Tf_{t})^{*}(s))\leq
\lambda^{t}((Tf^{t})^{*}(s))+\lambda_{t}((Tf_{t})^{*}(s))\leq s+s=2s.
\]
By the definition of $(Tf)^{*}$, therefore, we have
\[
\textup{(b)}\qquad
(Tf^{t})^{*}(s)+(Tf_{t})^{*}(s)\geq(Tf)^{*}(2s)\qquad\text{for all } s>0.
\]
Suppose $r_{1} < \infty$ and $q < \infty$. Then, from (b) with $s = t$, a
change of variables and Minkowski's inequality, we obtain
\begin{align*}
\|Tf\|_{pq}^{*}&=(q/p)^{1/q}\Biggl\{\int_{0}^{\infty}[t^{1/p}(Tf)^{*}(t)]^{q}\frac{dt}{t}\Biggr\}^{1/q}\\
&=(q/p)^{1/q}2^{1/p}\Biggl\{\int_{0}^{\infty}[t^{1/p}(Tf)^{*}(2t)]^{q}\frac{dt}{t}\Biggr\}^{1/q}\\
&\leq(q/p)^{1/q}2^{1/p}\Biggl\{\Biggl(\int_{0}^{\infty}[t^{1/p}(Tf^{t})^{*}(t)]^{q}\frac{dt}{t}\Biggr)^{1/q}\\
&\quad+\Biggl(\int_{0}^{\infty}[t^{1/p}(Tf_{t})^{*}(t)]^{q}\frac{dt}{t}\Biggr)^{1/q}\Biggr\}.
\end{align*}
Since $T$ is of restricted weak types $(r_{0}, p_{0})$ and $(r_{1}, p_{1})$ we
have $t^{1/p_{0}}(Tf^{t})^{*}(t) \leq k_{0}\|f^{t}\|_{r_{0}1}^{*}$ and
$t^{1/p_{1}}(Tf_{t})^{*}(t) \leq k_{1}\|f_{t}\|_{r_{1}1}^{*}$ for all $t > 0$.
Thus, the above sum within brackets is not larger than
\[
\begin{aligned}&k_{0}\left(\int_{0}^{\infty}[t^{(1/p)-(1/p_{0})}\|f^{t}\|_{r_{0}1}^{*}]^{q}\frac{dt}{t}\right)^{1/q}\\
&\qquad+k_{1}\left(\int_{0}^{\infty}[t^{(1/p)-(1/p_{1})}\|f_{t}\|_{r_{1}1}^{*}]^{q}\frac{dt}{t}\right)^{1/q}.\end{aligned}
\]
But, it follows from the first part of (a), a change of variables and
Lemma~\ref{lem:c5-3-14}, part (i), that
\begin{align*}
&\left(\int_{0}^{\infty}[t^{(1/p)-(1/p_{0})}\|f^{t}\|_{r_{0}1}^{*}]^{q}\frac{dt}{t}\right)^{1/q}\\
&\quad\leq\left(\int_{0}^{\infty}\left[t^{(1/p)-(1/p_{0})}\left(\frac{1}{r_{0}}\int_{0}^{t^{\gamma}}f^{*}(s)s^{(1/r_{0})-1}\,ds\right)\right]^{q}\frac{dt}{t}\right)^{1/q}\\
&\quad=|\gamma|^{-(1/q)}(1/r_{0})\left(\int_{0}^{\infty}\left[u^{(1/r)-(1/r_{0})}\int_{0}^{u}f^{*}(s)s^{(1/r_{0})-1}\,ds\right]^{q}\frac{du}{u}\right)^{1/q}\\
&\quad\leq\frac{|\gamma|^{-(1/q)}}{1-(r_{0}/r)}
\left(\int_{0}^{\infty}[u^{1/r_{0}}f^{*}(u)]^{q}u^{(q/r)-(q/r_{0})}\frac{du}{u}\right)^{1/q}\\
&\quad=\frac{1}{1-(r_{0}/r)}\left(\frac{r}{|\gamma|q}\right)^{1/q}\|f\|_{rq}^{*}.
\end{align*}
A similar argument, involving the second parts of (a) and
Lemma~\ref{lem:c5-3-14}, gives us
\begin{align*}
&\left(\int_{0}^{\infty}[t^{(1/p)-(1/p_{1})}\|f_{t}\|_{r_{1}1}^{*}]^{q}\frac{dt}{t}\right)^{1/q}\\
&\quad\leq\left(\int_{0}^{\infty}\left[t^{(1/p)-(1/p_{1})}\left(\frac{1}{r_{1}}\int_{0}^{t^{\gamma}}f^{*}(t^{\gamma})s^{(1/r_{1})-1}\,ds\right.\right.\right.\\
&\quad\quad\left.\left.\left.{}+\frac{1}{r_{1}}\int_{t^{\gamma}}^{\infty}f^{*}(s)s^{(1/r_{1})-1}\,ds\right)\right]^{q}\frac{dt}{t}\right)^{1/q}\\
&\quad\leq\left(\int_{0}^{\infty}[t^{(1/p)-(1/p_{1})}f^{*}(t^{\gamma})t^{\gamma/r_{1}}]^{q}\frac{dt}{t}\right)^{1/q}\\
&\quad\quad+\frac{1}{r_{1}}\left(\int_{0}^{\infty}\left[t^{(1/p)-(1/p_{1})}\int_{t^{\gamma}}^{\infty}f^{*}(s)s^{(1/r_{1})-1}\,ds\right]^{q}\frac{dt}{t}\right)^{1/q}\\
&\quad=|\gamma|^{-(1/q)}\Biggl\{\left(\int_{0}^{\infty}[u^{1/r}f^{*}(u)]^{q}\frac{du}{u}\right)^{1/q}\\
&\quad\quad+\frac{1}{r_{1}}\left(\int_{0}^{\infty}\left[u^{(1/r)-(1/r_{1})}\int_{u}^{\infty}f^{*}(s)s^{(1/r_{1})-1}\,ds\right]^{q}\frac{du}{u}\right)^{1/q}\Biggr\}\\
&\quad\leq\left(\frac{r}{|\gamma|q}\right)^{1/q}\|f\|_{rq}^{*}
+\frac{|\gamma|^{-(1/q)}}{(r_{1}/r)-1}
\!\left(\int_{0}^{\infty}[f^{*}(u)u^{1/r_{1}}]^{q}u^{(q/r)-(q/r_{1})}\frac{du}{u}\right)^{1/q}\\
&\quad=\left(\frac{r}{|\gamma|q}\right)^{1/q}\left(\frac{r_{1}}{r_{1}-r}\right)\|f\|_{rq}^{*}.
\end{align*}
We have shown, therefore, that
\[
\|Tf\|_{pq}^{*}\leq B\|f\|_{rq}^{*}
\]
with
\[
B=\left(\frac{r}{|\gamma|p}\right)^{1/q}2^{1/p}
\left(\frac{rk_{0}}{r-r_{0}}+\frac{r_{1}k_{1}}{r_{1}-r}\right).
\]

When $r_{1} < \infty$ and $q = \infty$ we estimate $t^{1/p}(Tf)^{*}(t)$, for
$t > 0$, using inequalities (a) and (b), just as was done in the above proof.
We obtain positive constants $c_{1}$, $c_{2}$, and $c_{3}$ (in fact,
$c_{1} = 2^{1/p}k_{0}$, and $c_{2} = c_{3} = 2^{1/p}k_{1}$) such that
\begin{align*}
t^{1/p}(Tf)^{*}(t)&\leq c_{1}t^{(1/p)-(1/p_{0})}\int_{0}^{t^{\gamma}}f^{*}(s)s^{(1/r_{0})-1}\,ds\\
&\quad+c_{2}t^{(1/p)-(1/p_{1})}\int_{0}^{t^{\gamma}}f^{*}(t^{\gamma})s^{(1/r_{1})-1}\,ds\\
&\quad+c_{3}t^{1/p-1/p_{1}}\int_{t^{\gamma}}^{\infty}f^{*}(s)s^{(1/r_{1})-1}\,ds.
\end{align*}
If we now use the fact that $f^{*}(s)s^{1/r} \leq \|f\|_{r\infty}^{*}$ we have
\[
\begin{aligned}&t^{1/p}(Tf)^{*}(t)\leq\left\{\frac{c_{1}}{(1/r_{0})-(1/r)}+c_{2}r_{1}
+\frac{c_{3}}{(1/r)-(1/r_{1})}\right\}\|f\|_{r\infty}^{*}\\&=B\|f\|_{r\infty}^{*},\end{aligned}
\]
and, thus, $\|Tf\|_{p\infty}^{*} \leq B\|f\|_{r\infty}^{*}$.

In the remaining case, $r_{1} = \infty$ and $q = \infty$, the proof follows the
same lines as that of the above cases, except for the necessity of using the
estimate $\|f_{t}\|_{r_{1}\infty}^{*} \leq f^{*}(t^{\gamma})$.
\end{proof}

Theorem~\ref{thm:c5-3-15} obviously contains the Marcinkiewicz interpolation
theorem as a special case; moreover, it gives us sharper results since it
involves the spaces $L(p, q)$. For example, if we apply it to the Fourier
transform we obtain the following strengthening of the inequality of
Hausdorff-Young (see \eqref{eq:c5-1-2}):

\begin{corollary}\label{cor:c5-3-16}
If $f \in L^{p}(E_{n})$, $1 < p \leq 2$, then $\hat{f}$ belongs to
$L(p', p)$ and there exists a constant $B = B_{p}$ such that
\[
\|\hat{f}\|_{p'p}^{*}\leq B\|f\|_{p},
\]
where $(1/p) + (1/p') = 1$.
\end{corollary}

The remainder of this section will be devoted to showing that most of the
spaces $L(p, q)$ have a norm that is closely connected to $\|\ \|_{pq}^{*}$;
moreover, with this norm, $L(p, q)$ is a Banach space. The following lemmas
will be needed to establish these facts.

For notational convenience we restrict our attention to nonatomic measure
spaces\index{non-atomic measure space@non-atomic measure space}. This means that $M$ does not contain atoms\index{atoms@atoms} (sets of positive measure
all of whose measurable subsets have measure zero). The applications we shall
make of the material developed in this section will involve only the nonatomic
measure spaces $(E_{n}, \mathcal{B}, m)$, where $\mathcal{B}$ consists of the
Borel subsets of $E_{n}$ and $m$ is Borel-Lebesgue measure. The only place
where this restriction is used is in the proof of part (iii) of the following
lemma.

\begin{lemma}\label{lem:c5-3-17}
If $E \in \mathcal{M}$ and
$F_{s} = \{x \in M : |f(x)| > s \geq 0\}$ (so that $\lambda(s) = \mu(F_{s})$)
then
\[
\textup{(i)}\qquad \int_{E}|f|\,d\mu\leq\int_{0}^{\mu(E)}f^{*}(t)\,dt;
\]
\[
\textup{(ii)}\qquad \int_{F_{s}}|f|\,d\mu=\int_{0}^{\lambda(s)}f^{*}(t)\,dt;
\]
\textup{(iii)} if $\mu(M) \geq t > 0$ then there exists a set
$E_{t} \in \mathcal{M}$ such that $\mu(E_{t}) = t$ and
\[
\int_{E_{t}}|f|\,d\mu=\int_{0}^{t}f^{*}(u)\,du.
\]
\end{lemma}

\begin{proof}
It is clear that if $|f_{1}| \leq |f|$ then $f_{1}^{*} \leq f^{*}$. Thus, if
$\chi_{E}$ is the characteristic function of $E$ and $f_{1} = f\chi_{E}$, then
\[
\textup{(a)}\qquad
\int_{0}^{\mu(E)}f_{1}^{*}(t)\,dt\leq\int_{0}^{\mu(E)}f^{*}(t)\,dt.
\]
On the other hand, the distribution function of $f_{1}$ must be bounded by
$\mu(E)$ and, consequently, $f_{1}^{*}(t) = 0$ for $t > \mu(E)$. Using
\eqref{eq:c5-3-9} with $p = 1$, therefore,
\[
\textup{(b)}\qquad
\int_{E}|f|\,d\mu=\int_{M}|f_{1}|\,d\mu=\int_{0}^{\infty}f_{1}^{*}(t)\,dt
=\int_{0}^{\mu(E)}f_{1}^{*}(t)\,dt.
\]
Inequality (a) and inequality (b) imply (i).

In order to show (ii) we first observe that if we put
\[
g(t)=\begin{cases}
f^{*}(t) & \quad\text{for } 0 < t < \lambda(s),\\
0 & \quad\text{for } \lambda(s) \leq t,\\
\end{cases}
\]
then $g$ and $h = f\chi_{F_{s}}$ have the same distribution function (when $f$
is simple this follows immediately from Lemma~\ref{lem:c5-3-4}, part (iii), and
equalities \eqref{eq:c5-3-6} and \eqref{eq:c5-3-7}. The general case follows
from this and Lemma~\ref{lem:c5-3-5}). Thus,
\[
\int_{F_{s}}|f|\,d\mu=\int_{M}|h|\,d\mu=\int_{0}^{\infty}g(u)\,du
=\int_{0}^{\lambda(s)}f^{*}(u)\,du.
\]
Finally, if $0 < t \leq \mu(M)$, let
$G = \{x \in M : |f(x)| > f^{*}(t)\}$ and
$H = \{x \in M : |f(x)| \geq f^{*}(t)\}$. We claim that
\[
\mu(G)\leq t\leq\mu(H).
\]
The first inequality is a consequence of part (ii) of Lemma~\ref{lem:c5-3-4}
and the fact that $\mu(G) = \lambda(f^{*}(t))$. The second inequality is
particularly easy to establish with the help of \eqref{eq:c5-3-6} and
\eqref{eq:c5-3-7} when $f$ is simple; the general case then follows from
Lemma~\ref{lem:c5-3-5}. Since we are assuming that the measure space
$(M, \mathcal{M}, \mu)$ is not atomic we can find a set $E_{t}$ in
$\mathcal{M}$ such that $G \subset E_{t} \subset H$ and $\mu(E_{t}) = t$. Let
$h = f\chi_{E_{t}}$. Then $h^{*} = f^{*}$ when these functions are restricted
to $(0, t)$, while $h(u) = 0$ for $u \geq t$. Therefore,
\[
\int_{E_{t}}|f|\,d\mu=\int_{M}|h|\,d\mu=\int_{0}^{\infty}h^{*}(u)\,du
=\int_{0}^{t}f^{*}(u)\,du.
\]
\end{proof}

An immediate corollary of parts (i) and (iii) of this lemma is
\begin{equation}\label{eq:c5-3-18}
\sup_{E \in \mathcal{M},\ \mu(E) \leq t}\int_{E}|f|\,d\mu=\int_{0}^{t}f^{*}(u)\,du.
\end{equation}
If $f$ is not zero almost everywhere and $t > 0$, then the expression on the
left in \eqref{eq:c5-3-18} is positive. It follows, therefore, that the mapping
$f \to t^{-1}\int_{0}^{t}f^{*}(u)\,du$ is a norm defined on each of the spaces
$L^{p}(M)$, $1 \leq p \leq \infty$ (the fact that it is finite-valued is a
consequence of H\"older's inequality):
\[
\int_{E}|f|\,d\mu=\int_{M}\chi_{E}|f|\,d\mu\leq\|f\|_{p}\|\chi_{E}\|_{p'}
\leq t^{1/p'}\|f\|_{p},
\]
where $(1/p) + (1/p') = 1$. This norm is closely related to the
Hardy-Littlewood maximal function introduced in Chapter II, Section 3. In
order to see this relation we make a few simple observations concerning the
maximal function.

The maximal function introduced in the second chapter was defined to be the
supremum over all integral averages of $|f|$ taken over spheres with center at
$x$. In the one-dimensional case it is not hard to show that its basic
properties, developed in the third section of that chapter, are shared by the
more general ``nonsymmetric'' maximal function defined by letting its value at
$x$ be
\begin{equation}\label{eq:c5-3-19}
\sup\frac{1}{h+k}\int_{x-h}^{x+k}|f(u)|\,du,
\end{equation}
where the supremum is taken over all nonnegative $h$ and $k$ satisfying
$h + k > 0$. This is a consequence of the inequalities
\begin{align*}
\frac{1}{h+k}\int_{x-h}^{x+k}|f(u)|\,du
&\leq\frac{1}{h+k}\int_{x-(h+k)}^{x+h+k}|f(u)|\,du\\
&\leq 2\sup_{r>0}\frac{1}{2r}\int_{|x-u| \leq r}|f(u)|\,du
\end{align*}
which show that this nonsymmetric maximal function is no larger than twice the
one we introduced earlier. In order to avoid introducing new notation, we let
$m = m_{f}$ denote the function defined by \eqref{eq:c5-3-19}.

Suppose $g$ is a nonnegative function defined on the real line such that
$g(u) = 0$ for $u < 0$ and $g$ is nonincreasing when restricted to
$[0, \infty)$. If $t > 0$ it then follows from the fact that $g(u) = 0$ for
$u < 0$ that $m(t) = m_{g}(t)$ is the supremum over all integral averages
\[
\frac{1}{h+k}\int_{t-h}^{t+k}g(u)\,du
\]
such that $t - h \geq 0$. Since $g$ is nonincreasing, however, a simple
argument shows that each such average is not larger than
\[
\frac{1}{t}\int_{0}^{t}g(u)\,du.
\]
If we apply these arguments to the function $g$ whose values at $u \geq 0$ are
$f^{*}(u)$ we obtain the fact that
\[
m_{g}(t)=m_{f^{*}}(t)=m(t)=\frac{1}{t}\int_{0}^{t}f^{*}(u)\,du
\]
for all $t > 0$.

It is clear from the nonincreasing property of $f^{*}$ that
$m(t) \geq f^{*}(t)$ for all $t > 0$. Thus, if we define
\[
\|f\|_{pq}=\left(\frac{q}{p}\int_{0}^{\infty}[t^{1/p}m(t)]^{q}\frac{dt}{t}\right)^{1/q}
\]
when $1 \leq p < \infty$ and $1 \leq q < \infty$, and
\[
\|f\|_{pq}=\sup_{t>0}t^{1/p}m(t)
\]
when $1 \leq p \leq \infty$ and $q = \infty$, we must have
\begin{equation}\label{eq:c5-3-20}
\|f\|_{pq}^{*}\leq\|f\|_{pq}.
\end{equation}

One of the advantages of using $\|\ \|_{pq}$ instead of $\|\ \|_{pq}^{*}$ is
that the former is a norm. This follows from the fact, observed immediately
following \eqref{eq:c5-3-18}, that the mapping $t \to m(t) = m_{f}(t)$ is a
norm for each $t > 0$. From \eqref{eq:c5-3-9}, with $p = 1$, we see that
$\|f\|_{1\infty} = \|f\|_{1}$; on the other hand, a simple calculation shows
that a function $f$ satisfying $\|f\|_{1q} < \infty$, for $1 \leq q < \infty$,
must be zero almost everywhere. When $1 < p \leq \infty$, however,
$\|\ \|_{pq}$ is ``equivalent'' to $\|\ \|_{pq}^{*}$. More precisely, we have
the following result\footnote{It is important to use both $\|\ \|_{pq}^{*}$
and $\|\ \|_{pq}$ since each has its advantages. The advantage that
$\|\ \|_{pq}$ is a norm is offset by the fact that the former is often easier
to handle; moreover, as was just observed, $\|\ \|_{1q}$ cannot be used to
define a satisfactory extension of $L^{1}$.}:

\begin{theorem}\label{thm:c5-3-21}
If $f \in L(p, q)$, $1 < p \leq \infty$, then
\[
\|f\|_{pq}^{*}\leq\|f\|_{pq}\leq\frac{p}{p-1}\|f\|_{pq}^{*}.
\]
\end{theorem}

\begin{proof}
The first inequality is just inequality \eqref{eq:c5-3-20} and has already
been proved. When $1 < p < \infty$ and $1 \leq q < \infty$ the second
inequality follows from Lemma~\ref{lem:c5-3-14}, part (i):
\begin{align*}
\|f\|_{pq}&=\left(\frac{q}{p}\int_{0}^{\infty}\left[\int_{0}^{t}f^{*}(u)\,du\right]^{q}
t^{-q(1-1/p)-1}\,dt\right)^{1/q}\\
&\leq\frac{p}{p-1}\left(\frac{q}{p}\int_{0}^{\infty}[uf^{*}(u)]^{q}
u^{-q(1-1/p)-1}\,du\right)^{1/q}
=\frac{p}{p-1}\|f\|_{pq}^{*}.
\end{align*}
For the remaining cases, $1 < p \leq \infty$ and $q = \infty$, we have
\begin{align*}
t^{1/p}m(t)&=t^{(1/p)-1}\int_{0}^{t}f^{*}(u)\,du
=t^{(1/p)-1}\int_{0}^{t}u^{-1/p}u^{1/p}f^{*}(u)\,du\\
&\leq\|f\|_{p\infty}^{*}t^{(1/p)-1}\int_{0}^{t}u^{-1/p}\,du
=\frac{p}{p-1}\|f\|_{p\infty}^{*}.
\end{align*}
\end{proof}

\begin{theorem}\label{thm:c5-3-22}
$L(p, q)$, $1 < p \leq \infty$, $1 \leq q \leq \infty$, with the norm
$\|\ \|_{pq}$ is a Banach space.
\end{theorem}

\begin{proof}
We need to show that $L(p, q)$ is complete. When $p = \infty$ we only defined
$L(p, q)$ for $q = \infty$; since $L(\infty, \infty) = L^{\infty}$ the theorem,
in this case, follows from the completeness of $L^{\infty}$. Suppose,
therefore, that $1 < p < \infty$ and that $\{f_{n}\}$ is a Cauchy sequence in
$L(p, q)$. By Theorem~\ref{thm:c5-3-21}, part (i) of
Theorem~\ref{thm:c5-3-11}, and Lemma~\ref{lem:c5-3-8} we must then have
\[
\sup_{s>0}s^{p}\mu(\{x \in M : |f_{n}(x)-f_{m}(x)| > s\}) \to 0
\qquad\text{as } n, m \to \infty.
\]
In particular, the sequence $\{f_{n}\}$ is fundamental in measure; therefore,
there exists a subsequence $\{f_{n_{k}}\}$ and a function $f$ such that
$f_{n_{k}}(x) \to f(x)$ a.e. Given $\delta > 0$ there exists
$j_{0} = j_{0}(\delta)$ such that $\|f_{n} - f_{j}\|_{pq} < \delta$ for
$j, n \geq j_{0}$. Let $g_{k} = f_{n_{k}} - f_{j}$ and $g = f - f_{j}$; it
follows from Fatou's lemma that
\[
\int_{E}|g|\,d\mu\leq\liminf_{k \to \infty}\int|g_{k}|\,d\mu
\]
for any measurable set $E$. From \eqref{eq:c5-3-18} and another application of
Fatou's lemma, therefore,
\[
\|f-f_{j}\|_{pq}=\|g\|_{pq}\leq\delta
\]
for $j \geq j_{0}$. This proves the theorem. It can be shown that $L(p, q)$
cannot be normed when $p = 1$, and $1 < q \leq \infty$. For this, see 5.12
below.
\end{proof}

\section{Interpolation of Analytic Families of Operators}

In this section we will extend the M. Riesz convexity theorem to the case
where the operators being interpolated vary with the indices $p$ and $q$ in a
sufficiently smooth way. We suppose that to each $z$ in the strip
$S = \{z \in \mathbf{C} : 0 \leq \operatorname{Re} z \leq 1\}$ there is
assigned a linear operator $T_{z}$ on the space of simple functions in
$L^{1}(M)$ into measurable functions on $N$ in such a way that $(T_{z}f)g$ is
integrable on $N$ whenever $f$ is a simple function in $L^{1}(M)$ and $g$ is a
simple function in $L^{1}(N)$. We shall say that the family $\{T_{z}\}$ is
admissible if the mapping
\[
z\to\int_{N}(T_{z}f)g\,d\nu
\]
is analytic in the interior of $S$, continuous on $S$ and there exists a
constant $a < \pi$ such that
\[
e^{-a|y|}\log\left|\int_{N}(T_{z}f)g\,d\nu\right|,
\]
$z = x + iy$, is uniformly bounded above in the strip $S$.

The extension of Theorem~\ref{thm:c5-1-3} is, then, the following one:

\begin{theorem}\label{thm:c5-4-1}
Suppose $\{T_{z}\}$, $z \in S$, is an admissible family of\index{analytic family of operators@analytic family of operators} linear operators
satisfying
\[
\|T_{iy}f\|_{q_{0}}\leq M_{0}(y)\|f\|_{p_{0}}
\qquad\text{and}\qquad
\|T_{1+iy}f\|_{q_{1}}\leq M_{1}(y)\|f\|_{p_{1}}
\]
for all simple functions $f$ in $L^{1}(M, m, \mu)$, where
$1 \leq p_{j}$, $q_{j} \leq \infty$, $M_{j}(y)$, $j = 0$, 1, are independent
of $f$ and satisfy
\[
\sup_{-\infty<y<\infty}e^{-b|y|}\log M_{j}(y)<\infty
\]
for some $b < \pi$. Then, if $0 \leq t \leq 1$, there exists a constant
$M_{t}$ such that
\[
\|T_{t}f\|_{q_{t}}\leq M_{t}\|f\|_{p_{t}}
\]
for all such simple functions $f$ provided
$1/p_{t} = (1 - t)/p_{0} + t/p_{1}$ and
$1/q_{t} = (1 - t)/q_{0} + t/q_{1}$.
\end{theorem}

The proof of this theorem is very similar to that of the M. Riesz convexity
theorem once we have established the following extension of the three lines
theorem (Lemma~\ref{lem:c5-1-4}):

\begin{lemma}\label{lem:c5-4-2}
Suppose $F$ is a continuous function on $S$ that is analytic in the interior
of $S$ satisfying
\begin{equation}\label{eq:c5-4-3}
\sup_{0 \leq x \leq 1,\ -\infty<y<\infty}e^{-a|y|}\log|F(x+iy)|<\infty
\end{equation}
for some $a < \pi$. Then
\begin{equation}\label{eq:c5-4-4}
\begin{aligned}
\log|F(x)|&\\
&\leq\tfrac{1}{2}\sin\pi x\int_{-\infty}^{\infty}
\left\{\frac{\log|F(iy)|}{\cosh\pi y-\cos\pi x}
+\frac{\log|F(1+iy)|}{\cosh\pi y+\cos\pi x}\right\}dy
\end{aligned}
\end{equation}
whenever $0 < x < 1$.
\end{lemma}

\begin{proof}
For $\zeta \in D = \{z \in \mathbb{C} : |z| \leq 1\}$, $\zeta \neq 1$, $-1$,
let
\[
h(\zeta)=(1/\pi i)\log i(1+\zeta)/(1-\zeta).
\]
The function $h$ is easily seen to be the composition of the conformal mapping
$\zeta \to w = i(1 + \zeta)/(1 - \zeta)$ of the closed unit disk $D$, minus the
points 1 and $-1$, onto the upper half plane
$\{w \in \mathbb{C} : w \neq 0$, $\operatorname{Im} w \geq 0\}$ and of the
conformal mapping $w \to z = (\log w)/\pi i$ of this half-plane onto $S$.
Thus, $h$ maps $D - \{1, -1\}$ conformally onto $S$. Moreover, for $z \in S$
\[
\zeta=h^{-1}(z)=\frac{e^{\pi iz}-i}{e^{\pi iz}+i}.
\]
If we define $G(\zeta) = F(h(\zeta))$ for $\zeta \in D - \{1, -1\}$ we obtain
an analytic function on the open unit disk that is continuous on the closed
unit disk with the exception of the points 1 and $-1$. If $0 \leq \rho < R < 1$,
an application of the Poisson-Jensen formula gives us, for
$\zeta = \rho e^{i\theta}$,
\begin{equation}\label{eq:c5-4-5}
\log|G(\zeta)|\leq\frac{1}{2\pi}\int_{-\pi}^{\pi}
\frac{R^{2}-\rho^{2}}{R^{2}-2R\rho\cos(\theta-\varphi)+\rho^{2}}
\log|G(Re^{i\varphi})|\,d\varphi.
\end{equation}
Condition \eqref{eq:c5-4-3}, expressed in terms of $G$ and
$\eta = h^{-1}(x + iy)$, implies
\[
\log|G(\eta)|\leq A\{|1+\eta|^{-(a/\pi)}+|1-\eta|^{-(a/\pi)}\}
\]
for some constant $A$, independent of $\eta \in D$. Since $a/\pi < 1$ this
inequality, with $\eta = Re^{i\varphi}$, allows us to apply the dominated
convergence theorem of Lebesgue to the integral in \eqref{eq:c5-4-5} and we
see, by letting $R$ tend to 1, that
\begin{equation}\label{eq:c5-4-6}
\log|G(\zeta)|\leq\frac{1}{2\pi}\int_{-\pi}^{\pi}
\frac{1-\rho^{2}}{1-2\rho\cos(\theta-\varphi)+\rho^{2}}
\log|G(e^{i\varphi})|\,d\varphi
\end{equation}
for $\zeta = \rho e^{i\theta}$, $\rho < 1$. We now claim that the desired
inequality \eqref{eq:c5-4-4} can be obtained from \eqref{eq:c5-4-6} by a change
of variables. First we note that the condition
$0 < x = h(\rho e^{i\theta}) < 1$ can be easily translated to a condition on
$\rho$ and $\theta$: since, in this case, we must have
\[
\rho e^{i\theta}=h^{-1}(x)=\frac{e^{\pi ix}-i}{e^{\pi ix}+i}
=-i\frac{\cos\pi x}{1+\sin\pi x}
=\left\{\frac{\cos\pi x}{1+\sin\pi x}\right\}e^{-i(\pi/2)},
\]
it follows that $\rho = (\cos\pi x)/(1 + \sin\pi x)$ and $\theta = -(\pi/2)$,
when $0 < x \leq \frac{1}{2}$, while
$\rho = -(\cos\pi x)/(1 + \sin\pi x)$ and $\theta = \pi/2$, when
$\frac{1}{2} \leq x < 1$. In the first case we have
\begin{align*}
\frac{1-\rho^{2}}{1-2\rho\cos(\theta-\varphi)+\rho^{2}}
&=\frac{1-\rho^{2}}{1-2\rho\cos(\varphi+\pi/2)+\rho^{2}}\\
&=\frac{1-\rho^{2}}{1+2\rho\sin\varphi+\rho^{2}}\\
&=\frac{\sin\pi x}{1+\cos\pi x\sin\varphi};
\end{align*}
moreover, this equality is also valid when $\frac{1}{2} \leq x < 1$. As we can
see from the relation
$e^{i\varphi} = h^{-1}(iy) = (e^{-\pi y} - i)/(e^{-\pi y} + i)$, as $\varphi$
ranges from $-\pi$ to 0, $y$ ranges from $+\infty$ to $-\infty$. Furthermore,
$\sin\varphi = -1/(\cosh\pi y)$ and
$d\varphi = -[\pi/(\cosh\pi y)]\,dy$. We therefore obtain
\begin{align*}
&\frac{1}{2\pi}\int_{-\pi}^{0}\frac{1-\rho^{2}}{1-2\rho\cos(\theta-\varphi)+\rho^{2}}\log|G(e^{i\varphi})|\,d\varphi\\
&\qquad=\frac{1}{2}\int_{-\infty}^{\infty}\frac{\sin\pi x}{\cosh\pi y-\cos\pi x}\log|F(iy)|\,dy.
\end{align*}
Similarly, as $\varphi$ ranges from 0 to $\pi$, $h(e^{i\varphi})$ describes the
points $1 + iy$ with $-\infty < y < \infty$ and we obtain
\begin{align*}
&\frac{1}{2\pi}\int_{0}^{\pi}\frac{1-\rho^{2}}{1-2\rho\cos(\theta-\varphi)+\rho^{2}}\log|G(e^{i\varphi})|\,d\varphi\\
&\qquad=\frac{1}{2}\int_{-\infty}^{\infty}\frac{\sin\pi x}{\cosh\pi y+\cos\pi x}\log|F(1+iy)|\,dy.
\end{align*}
This proves the lemma.
\end{proof}

Before showing how \eqref{lem:c5-4-2} can be used in order to prove
Theorem~\ref{thm:c5-4-1} we make some observations. First, we show that
Lemma~\ref{lem:c5-4-2} is, indeed, an extension of the three lines theorem. We
have
\begin{align*}
&\frac{1}{2}\int_{-\infty}^{\infty}\frac{\sin\pi x}{\cosh\pi y+\cos\pi x}\,dy\\
&\quad=\frac{1}{2}\int_{-\infty}^{\infty}
\frac{\tan(\pi/2)x}{(1+\tan^{2}(\pi/2)x\,\tanh^{2}(\pi/2)y)\cosh^{2}(\pi/2)y}\,dy\\
&=\frac{1}{\pi}\int_{-1}^{1}\frac{\tan(\pi/2)x}{1+s^{2}\tan^{2}(\pi/2)x}\,ds
=\frac{2}{\pi}\int_{0}^{1}\frac{\tan(\pi/2)x}{1+s^{2}\tan^{2}(\pi/2)x}\,ds\\
&=\frac{2}{\pi}\int_{0}^{\tan(\pi/2)x}\frac{du}{1+u^{2}}
=\frac{2}{\pi}\tan^{-1}\bigl(\tan(\pi/2)x\bigr)=x.
\end{align*}
Since $\sin\pi(1-x) = \sin\pi x$ and $\cos\pi(1-x) = -\cos\pi x$, we also have
\[
\frac{1}{2}\int_{-\infty}^{\infty}\frac{\sin\pi x}{\cosh\pi y-\cos\pi x}\,dy=1-x.
\]
Consequently, if condition \eqref{eq:c5-4-3} were replaced by the conditions
$|F(iy)| \leq m_{0}$, $|F(1 + iy)| \leq m_{1}$ and the fact that $|F|$ is
uniformly bounded in $S$, then the equalities we have just established and
\eqref{eq:c5-4-4} imply $|F(x + iy)| \leq m_{0}^{1-x}m_{1}^{x}$. That is, the
three lines theorem (Lemma~\ref{lem:c5-1-4}) is a special case of
Lemma~\ref{lem:c5-4-2}.

It is clear from the proof of \eqref{lem:c5-4-2} that the integral on the right
side of inequality \eqref{eq:c5-4-4} is the ``conformal image'' of a
corresponding Poisson integral on the unit circle. In fact, a small change in
the argument used to establish \eqref{lem:c5-4-2} gives us the following
solution to the Dirichlet problem for the strip: If $f$ is defined on the
vertical lines $x = 0$ and $x = 1$ in such a way that $e^{-a|y|}f(y)$ and
$e^{-a|y|}f(1 + iy)$ are bounded for some $a < \pi$, then
\[
u(x,y)=\tfrac{1}{2}\sin\pi x\int_{-\infty}^{\infty}
\left\{\frac{f(iy)}{\cosh\pi y-\cos\pi x}
+\frac{f(1+iy)}{\cosh\pi y+\cos\pi x}\right\}dy
\]
defines a harmonic function in the interior of $S$ having (nontangential)
boundary values $f$. The conditions $e^{-a|y|}f(y) = O(1)$ and
$e^{-a|y|}f(1 + iy) = O(1)$ imply the existence of the above integral (this
also explains the role played by condition \eqref{eq:c5-4-3}).

We now pass to the proof of Theorem~\ref{thm:c5-4-1}. Let $f \in L^{1}(M)$ and
$g \in L^{1}(N)$ be simple functions satisfying
\[
\|f\|_{p}=1=\|g\|_{q'},
\]
where $p = p_{t}$, $q = q_{t}$ and $1/q + 1/q' = 1$. Suppose that
$f = \sum_{j=1}^{m}a_{j}\chi_{E_{j}}$ and
$g = \sum_{k=1}^{n}b_{k}\chi_{F_{k}}$. Using the same notation that was
introduced in the proof of the M. Riesz convexity theorem
(Theorem~\ref{thm:c5-1-3}), we construct the function
\[
F(z)=\int_{N}(T_{z}f_{z})g_{z}\,d\nu
=\sum_{j,k=1}^{m,n}|a_{j}|^{\alpha(z)/\alpha}|b_{k}|^{(1-\beta(z))/(1-\beta)}\gamma_{jk}(z),
\]
where
$\gamma_{jk}(z) = e^{i(\theta_{j} + \varphi_{k})}\int_{N}(T_{z}\chi_{E_{j}})\chi_{F_{k}}\,d\nu$.
It follows from the assumption that $\{T_{z}\}$ is an admissible family that
$F$ satisfies the hypotheses of Lemma~\ref{lem:c5-4-2}. Moreover, since
\[
|f_{iy}|^{p_{0}}=|f|^{p}=|f_{1+iy}|^{p_{1}}
\qquad
|g_{iy}|^{q_{0}'}=|g|^{q'}=|g_{1+iy}|^{q_{1}'}
\]
an application of H\"older's inequality gives us $|F(iy)| \leq M_{0}(y)$ and
$|F(1 + iy)| \leq M_{1}(y)$. Thus, by Lemma~\ref{lem:c5-4-2},
\begin{align*}
&\bigg|\int_{N}(T_{t}f)g\,d\nu\bigg|=|F(t)|\\
&\quad\leq\exp\left[\tfrac{1}{2}\sin\pi t\int_{-\infty}^{\infty}
\left\{\frac{\log M_{0}(y)}{\cosh\pi y-\cos\pi t}\right.\right.\\
&\quad\quad\left.\left.{}+\frac{\log M_{1}(y)}{\cosh\pi y+\cos\pi t}\right\}dy\right]=M_{t}.
\end{align*}
Since
$\|T_{t}f\|_{q_{t}} = \sup_{\|g\|_{q_{t}'} = 1}|\int_{N}(T_{t}f)g\,d\nu|$,
this proves the theorem.

\section{Further Results}

\textbf{5.1.} The theory of interpolation of linear operators has been
extended in many ways to general Banach spaces. The setting for these
extensions can be described in the following manner: Let $V$ be a topological
vector space and $A^{0}$, $A^{1}$ two Banach spaces, with norms $\|\ \|_{0}$
and $\|\ \|_{1}$, that are continuously embedded in $V$. We can convert
$A^{0} \cap A^{1}$ and
$A^{0} + A^{1} = \{a = a_{0} + a_{1} : a_{0} \in A^{0} \text{ and } a_{1} \in A^{1}\}$
into Banach spaces by introducing the norm
$\|a\|_{A^{0} \cap A^{1}} = \max\{\|a\|_{0}, \|a\|_{1}\}$ in the former space
and
$\|a\|_{A^{0} + A^{1}} = \inf\{\|a_{0}\|_{0} + \|a_{1}\|_{1} : a_{0} \in A^{0}, a_{1} \in A^{1} \ni a_{0} + a_{1} = a\}$
in the latter. An intermediate space\index{interpolation@interpolation!intermediate space@intermediate space} between $A^{0}$ and $A^{1}$ is any Banach
space $A$ satisfying $A^{0} \cap A^{1} \subset A \subset A^{0} + A^{1}$. For
example, when $A^{0} = L^{p_{0}}$, $A^{1} = L^{p_{1}}$,
$1 \leq p_{0}, p_{1} \leq \infty$,
$1/p_{t} = (1 - t)/p_{0} + t/p_{1}$, $0 \leq t \leq 1$, then $L^{p_{t}}$ is an
intermediate space between $L^{p_{0}}$ and $L^{p_{1}}$. Suppose an intermediate
space $A$ is invariant under each linear transformation $T$ of $A^{0} + A^{1}$
into itself having the property that its restriction to $A^{j}$, $j = 0$, 1, is
a bounded operator of $A^{j}$ into itself. If, moreover, the restriction of any
such $T$ to $A$ is also a bounded operator, we say that $A$ is a space of
linear interpolation\index{space of linear interpolation@space of linear interpolation} between $A^{0}$ and $A^{1}$. It follows from the M. Riesz
convexity theorem that $L^{p_{t}}$ is a space of linear interpolation between
$L^{p_{0}}$ and $L^{p_{1}}$.

We can then formulate the general problems in the abstract theory of
interpolation of linear operators in the following way:

\textup{(i)} Given two Banach spaces $A^{0}$ and $A^{1}$, how can ``all'' the
spaces of linear interpolation between $A^{0}$ and $A^{1}$ be characterized?

\textup{(ii)} How can one construct spaces of linear interpolation between two
given Banach spaces $A^{0}$ and $A^{1}$?

\textup{(iii)} Suppose $A$ is a space of linear interpolation between $A^{0}$
and $A^{1}$. If $B^{0}$ and $B^{1}$ are two other Banach spaces, does there
exist a space of linear interpolation $B$ (between $B^{0}$ and $B^{1}$) such
that each linear transformation of $A^{0} + A^{1}$ into $B^{0} + B^{1}$ whose
restriction to $A^{j}$ maps this space continuously into $B^{j}$, $j = 0$, 1,
has the property that it maps $A$ continuously into $B$?

\textup{(iv)} Suppose $A$ is a space of linear interpolation that has been
constructed by some specific method from $A^{0}$ and $A^{1}$ and the ``same''
method has been employed to obtain $B$ from two other Banach spaces $B^{0}$ and
$B^{1}$. Is it true that each linear transformation from $A^{0} + A^{1}$ into
$B^{0} + B^{1}$ that maps $A^{j}$ continuously into $B^{j}$, $j = 0$, 1, also
maps $A$ continuously into $B$?

An answer to the first problem has been obtained by Gagliardo (see
\cite{Gagliardo1}). Moreover, Gagliardo's method enabled him to obtain an
affirmative answer to problem (iii). These results are very general and are not
as closely related to the material developed in this chapter as are two
constructive theories of interpolation (that is, answers to problems (ii) and
(iv)) developed by Lions \cite{Lions1}, Calder\'on \cite{Calderon2}, Lions and
Peetre \cite{LionsPeetre1}. The next few paragraphs will be devoted to a
description of these methods of interpolation.

\textbf{5.2.} We begin by describing the complex method\index{interpolation@interpolation!complex method@complex method} of interpolation due
to Calder\'on. Let $B$ be a Banach space and $D$ a domain in the complex
plane. A mapping $z \to b(z)$, from $D$ into $B$, is said to be analytic if
$z \to l[b(z)]$ is an analytic complex-valued function on $D$ for each
continuous linear functional $l$ on $B$. We introduce an auxiliary space
$\mathcal{F}(A^{0}, A^{1})$ consisting of all functions $f$ defined on the
strip $S = \{z \in \mathbb{C} : 0 \leq \operatorname{Re} z \leq 1\}$ with
values in $A^{0} + A^{1}$ satisfying:

\begin{enumerate}[label=(\arabic*),leftmargin=*,itemsep=0pt]
\item $f$ is analytic in the interior
$S^{0} = \{z \in \mathbb{C} : 0 < \operatorname{Re} z < 1\}$ of $S$;
\item $f$ is continuous and $\|f\|_{A^{0}+A^{1}}$ is bounded on $S$;
\item $f(it) \in A^{0}$ for $-\infty < t < \infty$,
$t \to f(it)$ is continuous as a function from $(-\infty, \infty)$ into $A^{0}$
and $t \to \|f(it)\|_{0}$ is bounded;
\item $f(1 + it) \in A^{1}$ for $-\infty < t < \infty$,
$t \to f(1 + it)$ is continuous as a function from $(-\infty, \infty)$ into
$A^{1}$ and $t \to \|f(1 + it)\|_{1}$ is bounded.
\end{enumerate}

$\mathcal{F}(A^{0}, A^{1})$ is a Banach space if we introduce the norm
$\|f\|_{\mathcal{F}} = \|f\|_{\mathcal{F}(A^{0}, A^{1})} = \max\{\sup_{t}\|f(it)\|_{0}, \sup_{t}\|f(1 + it)\|_{1}\}$.
Suppose $0 \leq t \leq 1$ and $N_{t} = \{f \in \mathcal{F} : f(t) = 0\}$ we
then define $A_{t}$ to be the space $\mathcal{F}/N_{t}$. Clearly, we can then
identify $A_{t}$ with the linear space of all elements
$f(t) \in A^{0} + A^{1}$ with $f \in \mathcal{F}$. We obtain a Banach space if
we introduce on $A_{t}$ the usual quotient space norm
$\|a\|_{A_{t}} = \inf\{\|f\|_{\mathcal{F}} : f \in \mathcal{F}, f(t) = a\}$
whenever $a \in A_{t}$. We shall assume that $\{0\} \neq A^{0} \cap A^{1}$.

It is clear that $A_{t}$ is an intermediate space between $A^{0}$ and $A^{1}$.
It can be shown, moreover, that $A_{t}$ is a space of linear interpolation
between $A^{0}$ and $A^{1}$. In fact, we obtain the following solution to
problem (iv) announced above:

If $B^{0}$ and $B^{1}$ are two Banach spaces and $T$ is a linear
transformation from $A^{0} + A^{1}$ into $B^{0} + B^{1}$ that maps $A^{j}$ into
$B^{j}$, $j = 0$, 1, in such a way that $\|Ta\|_{j} \leq M_{j}\|a\|_{j}$ for
all $a \in A^{j}$ then $T$ maps $A_{t}$ into $B_{t}$ and
$\|Ta\|_{B_{t}} \leq M_{0}^{(1-t)}M_{1}^{t}\|a\|_{A_{t}}$ for all
$a \in A_{t}$.

Calder\'on gives another (but related) construction of spaces of linear
interpolation between $A^{0}$ and $A^{1}$ that leads to another solution to
problem (iv). When we apply either of these results, however, we must identify
the spaces of linear interpolation that have been obtained. In order to better
understand the significance of this new problem we present some basic examples
of interpolation of linear operators (for the most part, these were known
before the formulation of the general theory).

\textbf{5.3.} Suppose $(M, \mathcal{M}, \mu)$ is a measure space and
$A^{j} = L^{p_{j}}(M)$, $1 \leq p_{j} \leq \infty$ for $j = 0$, 1. We can
embed $A^{0}$ and $A^{1}$ in, say, the space $V$ of locally integrable
functions. Then, there exists a norm preserving isomorphism between $A_{t}$ and
$L^{p_{t}}(M)$, where $1/p_{t} = (1 - t)/p_{0} + t/p_{1}$, $0 \leq t \leq 1$,
provided $p_{t} < \infty$. In case $p_{t} = \infty$, $p_{t}$ must be either
$p_{0}$ or $p_{1}$; let us suppose $p_{t} = p_{0}$ and $p_{1} < \infty$. We do
not then obtain, in general, a norm preserving isomorphism between $A_{0}$ and
$A^{0} = L^{\infty}(M)$. It can be shown, however, that $A_{0}$ can be
identified with the closure in the $L^{\infty}$ norm of the set of functions
$f \in L^{\infty}(M)$ for which
$\mu(\{x \in M : |f(x)| > \alpha\}) < \infty$ for all $\alpha > 0$. If
$B^{0} = L^{q_{0}}(N)$ and $B^{1} = L^{q_{1}}(N)$, where
$(N, \mathcal{N}, \nu)$ is another measure space, the theorem stated in
paragraph 5.2 is a reformulation of the M. Riesz convexity theorem
(Theorem~\ref{thm:c5-1-3}).

\textbf{5.4.} Suppose $(M, \mathcal{M}, \mu)$ is a measure space and
$\mu_{j}$, $j = 0$, 1, is a measure on $\mathcal{M}$ that is absolutely
continuous with respect to $\mu$. If $A^{j} = L^{p_{j}}(M, \mu_{j})$ then
there exists a norm-preserving isomorphism between $A_{t}$ and
$L^{p_{t}}(M, \mu_{t})$ where
$0 < 1/p_{t} = (1 - t)/p_{0} + t/p_{1}$ and
$d\mu_{t} = \alpha_{0}^{1-t}\alpha_{1}^{t}\,d\mu$, with $0 \leq t \leq 1$
($\alpha_{j}$ being the Radon-Nikodym derivative of $\mu_{j}$, $j = 0$, 1, with
respect to $\mu$). When $p_{t} = \infty$ we have to make the modifications
described in 5.3. See Stein and G. Weiss \cite{SteinWeiss5}.

\textbf{5.5.} For $p \geq 1$ let $H^{p}$ be the space of all functions
$F(z)$ that are analytic in the upper half-plane $\operatorname{Im} z > 0$ and
satisfy
\[
\|F\|_{p}=\sup_{y>0}\left(\int_{-\infty}^{\infty}|F(x+iy)|^{p}\,dx\right)^{1/p}<\infty.
\]
This is the $H^{p}$ space associated with the tube whose base is the positive
real line (see Chapter III). If $A^{j} = H^{p_{j}}$, $1 \leq p_{j} < \infty$,
then $A_{t}$, $0 \leq t \leq 1$, is equivalent to $H^{p_{t}}$ (that is, there
exists an invertible linear map of $A_{t}$ onto $H^{p_{t}}$). (See Salem and
Zygmund \cite{SalemZygmund1}, Calder\'on and Zygmund \cite{CalderonZygmund2}
and G. Weiss \cite{Weiss1}.) The identification of the spaces of linear
interpolation between two $H^{p}$ spaces associated with higher dimensional
tubes is an open problem. This also applies to the $H^{p}$ spaces of systems of
conjugate harmonic functions that are discussed in the next chapter.

\textbf{5.6.} Suppose $0 < \alpha < 2$, we then consider the space
$\lambda_{\alpha} = \lambda_{\alpha}(E_{n})$ consisting of all continuous
bounded functions on $E_{n}$ satisfying
$\sup_{x}|f(x + t) + f(x - t) - 2f(x)| = o(|t|^{\alpha})$ as $|t| \to 0$. This
space is a Banach space with the norm
$\|f\|^{(\alpha)} = \|f\|_{\infty} + \sup_{x,t}\{|t|^{-\alpha}|f(x + t) + f(x - t) - 2f(x)|\}$.
For general $\alpha > 0$ we define $\lambda_{\alpha}$ to be the space of all
$C^{k}$ functions, where $k$ is the largest integer strictly smaller than
$\alpha$, satisfying $D^{\beta}f \in \lambda_{\alpha-k}$, for
$0 \leq |\beta| = \beta_{1} + \cdots + \beta_{n} < \alpha$ (we are using the
notation introduced in \eqref{eq:c1-1-9} of Chapter I). Then $\lambda_{\alpha}$
is a Banach space with the norm
$\|f\|^{(\alpha)} = \sum_{|\beta| < \alpha}\|D^{\beta}f\|^{(\alpha-k)}$. If
$A^{j} = \lambda_{\alpha_{j}}$, $j = 0$, 1, then $A_{t}$ is equivalent to
$\lambda_{\alpha}$, where $\alpha = (1 - t)\alpha_{0} + t\alpha_{1}$. (See
Taibleson \cite{Taibleson1} and Calder\'on \cite{Calderon2}.)

\textbf{5.7.} The following example is usually studied in the context of
$E_{n}$ but is simpler to describe in the setting of multiple Fourier series.
For any $\alpha \geq 0$ we define the space $L_{\alpha}^{p}(T_{n})$ to be the
set of those $f \in L^{p}(T_{n})$ with
$f(x) \sim \sum_{m}a_{m}e^{2\pi im \cdot x}$ so that
$a_{0} + \sum_{m \neq 0}|m|^{\alpha}a_{m}e^{2\pi im \cdot x}$ is the Fourier
series of a function $f^{(\alpha)} \in L^{p}(T_{n})$ (we are using the
notation that is introduced at the beginning of Chapter VII). This space is a
Banach space with the norm $\|f\|_{p}^{(\alpha)} = \|f^{(\alpha)}\|_{p}$.
Suppose $A^{j} = L_{\alpha_{j}}^{p}$, $1 < p_{j} < \infty$, $j = 0$, 1, then
$A_{t}$ is equivalent to $L_{\alpha}^{p}$, where
$1/p = (1 - t)/p_{0} + t/p_{1}$ and $\alpha = (1 - t)\alpha_{0} + t\alpha_{1}$.
The proof of this fact makes use of Theorem~\ref{thm:c5-4-1} and the fact that
$f \to \sum_{m}|m|^{-i\gamma}a_{m}e^{2\pi im \cdot x}$ is a bounded operator on
$L^{p}(T^{n})$ for $1 < p < \infty$ when $\gamma$ is real.

\textbf{5.8.} The classical interpolation theorem of M. Riesz, as formulated
in 5.3 or in Section 1 of this chapter, can be extended to the case of
Banach-space-valued functions on a measure space. For a detailed statement of
this result, see Calder\'on \cite{Calderon2}. Special cases, where the values
of the functions are themselves $L^{q}$ spaces, were obtained earlier by Boas
and Bochner \cite{BoasBochner1} and by Benedeck and Panzone
\cite{BenedekPanzone1}. Another extension of the classical interpolation
theorem of M. Riesz can be formulated for Banach Lattices\index{Banach lattice@Banach lattice} (these being Banach
function spaces such that $|f| \leq |g|$ implies $\|f\| \leq \|g\|$. For a
precise statement see Calder\'on \cite{Calderon2}). This extension includes the
M. Riesz theorem, the example described in 5.4 and an interpolation theorem
involving $L(p, q)$ spaces that were defined in the third section.

\textbf{5.9.} The constructive theory of interpolation mentioned in 5.1 which
was introduced by Gagliardo, Lions and Peetre can be described in the following
way. In terms of the notation introduced in 5.1 we define a norm on
$A^{0} + A^{1}$ for each $t > 0$ by letting
$K_{a}(t) = \inf\{\|a_{0}\|_{0} + t\|a_{1}\|_{1} : a_{0} + a_{1} = a, a_{0} \in A^{0}, a_{1} \in A^{1}\}$.
It follows immediately from this definition that $K_{a}$ is a concave and
nonnegative function on the positive real numbers. Suppose now that $\Phi$ is
a nonnegative (possibly infinite-valued) function defined on the set of all
nonnegative Lebesgue measurable functions on $(0, \infty)$ satisfying:
\[
\begin{aligned}
&\textup{(i)}\quad \Phi(h)=0 \text{ if and only if } h(t)=0 \text{ a.e.};\\
&\textup{(ii)}\quad \text{if } \Phi(h)<\infty \text{ then } h(t)<\infty \text{ a.e.};\\
&\textup{(iii)}\quad \Phi(ah)=a\Phi(h) \text{ when } a>0;\\
&\textup{(iv)}\quad \text{if } h(t)\leq\sum_{j=1}^{\infty}h_{j}(t) \text{ a.e. then } \Phi(h)\leq\sum_{j=1}^{\infty}\Phi(h_{j}).
\end{aligned}
\]
Such a function $\Phi$ is called a function norm\index{function norm@function norm} and the set
$A(\Phi) = \{a \in A^{0} + A^{1} : \Phi(K_{a}) < \infty\}$ is a subspace of
$A^{0} + A^{1}$. The mapping $a \to \Phi(K_{a})$ is a norm on $A(\Phi)$ which
shall be denoted by $\|a\|_{\Phi}$. When $\Phi$ satisfies certain general
conditions the space $A(\Phi)$ is a space of linear interpolation between
$A^{0}$ and $A^{1}$; moreover, we then have an interpolation theorem that gives
us a solution to problem (iv) that was posed in 5.1: Suppose $T$ is a linear
transformation from $A^{0} + A^{1}$ into $B^{0} + B^{1}$ that maps $A^{j}$ into
$B^{j}$, $j = 0$, 1, in such a way that $\|Ta\|_{j} \leq M_{j}$ for all
$a \in A^{j}$; then $T$ maps $A(\Phi)$ into $B(\Phi)$ and there exists a
constant $M = M(\Phi, M_{0}, M_{1})$ such that $\|Ta\|_{B(\Phi)} \leq M\|a\|_{A(\Phi)}$
for all $a \in A(\Phi)$. The proof of this theorem is particularly simple if
$\Phi$ satisfies the following condition: There exists a nonnegative function
$\psi$ on $(0, \infty)$ such that $\Phi(h^{\lambda}) \leq \psi(\lambda)\Phi(h)$,
where $\lambda > 0$ and $h^{\lambda}(t) = h(\lambda t)$. When this is the case,
we have $\|Ta\|_{B(\Phi)} = \Phi(K_{Ta}) \leq M_{0}\Phi(K_{a}^{M_{1}/M_{0}})
\leq M_{0}\psi(M_{1}/M_{0})\Phi(K_{a}) = M_{0}\psi(M_{1}/M_{0})\|a\|_{A(\Phi)}$
(the first inequality being a consequence of the hypotheses on $T$ and the
second following from our new assumption on $\Phi$). Thus, the interpolation
theorem is obtained with $M \leq M_{0}\psi(M_{1}/M_{0})$. If, for example,
$\psi(\lambda) = \lambda^{s}$ for $0 \leq s \leq 1$ we then have the familiar
inequality $M \leq M_{0}^{1-s}M_{1}^{s}$.

\textbf{5.10.} When $A^{0} = L^{1}$, $A^{1} = L^{\infty}$ and
$f \in L^{1} + L^{\infty}$ it is not hard to establish that
$K_{f}(t) = \int_{0}^{t}f^{*}(u)\,du$. The following observations show that it
suffices to consider only the case $f \geq 0$ and, in this case,
$K_{f}(t) = \inf\{\|f_{0}\|_{L^{1}} + t\|f_{1}\|_{L^{\infty}} : f = f_{0} + f_{1}, f_{j} \geq 0 \text{ and } f_{j} \in A^{j}, j = 0, 1\}$.
Suppose $f = f_{0} + f_{1}$; then
$|f| = (\operatorname{sgn}\bar{f})f = (\operatorname{sgn}\bar{f})f_{0} + (\operatorname{sgn}\bar{f})f_{1} = g_{0} + g_{1}$;
similarly, if $|f| = g_{0} + g_{1}$ then
$f = (\operatorname{sgn}\bar{f})g_{0} + (\operatorname{sgn}\bar{f})g_{1} = f_{0} + f_{1}$.
In either case we have
$\|f_{0}\|_{L^{1}} + t\|f_{1}\|_{L^{\infty}} = \|g_{0}\|_{L^{1}} + t\|g_{1}\|_{L^{\infty}}$
and, thus, $K_{f} = K_{|f|}$. When $f$ is real-valued and $f = f_{0} + f_{1}$
then $f = \operatorname{Re}f_{0} + \operatorname{Re}f_{1}$; thus, in estimating
$K_{f}$ it suffices to consider only real-valued $f_{0}$ and $f_{1}$. Finally,
if $f \geq 0$ and $f = f_{0} + f_{1}$ let $f_{0}' = \inf\{f_{0}^{+}, f\}$ and
$f_{1}' = f - f_{0}'$. Then, $f = f_{0}' + f_{1}'$,
$0 \leq f_{0}' \leq f_{0}^{+} \leq |f_{0}|$ and
$0 \leq f_{1}' = f - f_{0}' \leq |f - f_{0}| = |f_{1}|$. Consequently,
$\|f_{0}'\|_{L^{1}} + t\|f_{1}'\|_{L^{\infty}} \leq \|f_{0}\|_{L^{1}} + t\|f_{1}\|_{L^{\infty}}$.
We can, therefore, assume that $f = f_{0} + f_{1}$ with all three functions
being nonnegative. If $s = \|f_{1}\|_{L^{\infty}}$ define
$f^{(s)} = f \wedge s = \inf\{f, s\}$ and $g = f - f^{(s)}$. By changing $f_{1}$
on a set of measure 0, if necessary, we can assume $0 \leq f_{1}(x) \leq s$ for
all $x$. Since $f_{1} \leq f$ it follows that $f_{1} \leq f^{(s)}$ and,
consequently, $g = f - f^{(s)} \leq f - f_{1} = f_{0}$. This shows that
$\|g\|_{L^{1}} + t\|f^{(s)}\|_{L^{\infty}} = \|g\|_{L^{1}} + st \leq \|f_{0}\|_{L^{1}} + st \leq \|f_{0}\|_{L^{1}} + t\|f_{1}\|_{L^{\infty}}$.
Therefore, for $f \geq 0$ we have
$K_{f}(t) = \inf_{s \geq 0}\{\|f - f^{(s)}\|_{L^{1}} + ts : f^{(s)} = f \wedge s\}$.
If $\lambda$ denotes the distribution function of $f$ and
$s_{0} = \inf\{s : \lambda(s) < t\}$ we will now show that
$K_{f}(t) = \|f - f^{(s_{0})}\|_{L^{1}} + ts_{0}$. First suppose $s_{1} > s_{0}$;
then $f^{(s_{1})}(x) - f^{(s_{0})}(x)$ equals 0 when $f(x) \leq s_{0}$, equals
$f(x) - s_{0}$ if $s_{0} < f(x) \leq s_{1}$, and takes on the value $s_{1} - s_{0}$
if $s_{1} < f(x)$. Thus, since $\lambda(s_{0}) \leq t$,
\begin{align*}
\|f&-f^{(s_{1})}\|_{L^{1}}+ts_{1}-(\|f-f^{(s_{0})}\|_{L^{1}}+ts_{0})\\
&=t(s_{1}-s_{0})-\int(f^{(s_{1})}-f^{(s_{0})})\,d\mu\\
&\geq\lambda(s_{0})(s_{1}-s_{0})-\int_{f(x)>s_{0}}(f^{(s_{1})}-f^{(s_{0})})\,d\mu\\
&\geq\lambda(s_{0})(s_{1}-s_{0})-\lambda(s_{0})(s_{1}-s_{0})=0.
\end{align*}
Now suppose $s_{1} < s_{0}$ and $s_{1} \leq s < s_{0}$. For $s$ in this
interval $\lambda(s) \geq t$; therefore,
\begin{align*}
\|f&-f^{(s_{1})}\|_{L^{1}}+ts_{1}-(\|f-f^{(s)}\|_{L^{1}}+st)\\
&\geq\int_{f(x)>s_{1}}(f^{(s)}-f^{(s_{1})})\,d\mu-\lambda(s)(s-s_{1})\\
&\geq\int_{f(x)>s}(s-s_{1})\,d\mu-\lambda(s)(s-s_{1})\\
&=0.
\end{align*}
By continuity the inequality
$\|f - f^{(s_{1})}\|_{L^{1}} + ts_{1} \geq \|f - f^{(s)}\|_{L^{1}} + ts$ must
hold for $s = s_{0}$ and we have shown that
$K_{f}(t) = \|f - f^{(s_{0})}\|_{L^{1}} + ts_{0}$. It follows from the
definition of $f^{*}$ that $f^{*}(u) = s_{0}$ when $\lambda(s_{0}) \leq u < t$.
This, together with equality (ii) of Lemma~\ref{lem:c5-3-17}, gives us
\begin{align*}
K_{f}(t)&=\int(f-f^{(s_{0})})\,d\mu+ts_{0}
=\int_{f(x)>s_{0}}f\,d\mu-s_{0}\lambda(s_{0})+ts_{0}\\
&=\int_{0}^{\lambda(s_{0})}f^{*}(u)\,du+\int_{\lambda(s_{0})}^{t}f^{*}(u)\,du
=\int_{0}^{t}f^{*}(u)\,du.
\end{align*}
If we now choose the function norm
\[
\Phi(h)=\left((q/p)\int_{0}^{\infty}(t^{(1/p)-1}h(t))^{q}\,(dt/t)\right)^{1/q}
\]
we see that the space $A(\Phi)$ defined in 5.9 is $L(p, q)$, with
$1 < p \leq \infty$, and the norm is identical with $\|\ \|_{pq}$.

\textbf{5.11.} In the first section of this chapter we gave a geometrical
interpretation of the M. Riesz convexity theorem in terms of the points of the
unit square $Q = \{(\alpha, \beta) : 0 \leq \alpha, \beta \leq 1\}$. A similar
interpretation can be given for the Marcinkiewicz theorem
(Theorem~\ref{thm:c5-2-4}); however, in this case only the points of the
``lower triangle'' $\{(\alpha, \beta) \in Q : \alpha \geq \beta\}$ are
involved. One can show by example that this theorem is not valid in the
``upper triangle'' (see Hunt \cite{Hunt1}). In his paper \cite{Riesz2} M. Riesz
showed that his convexity theorem is valid for the entire unit square $Q$ only
if the $L^{p}$ spaces under consideration consist of complex-valued functions.
In case only real-valued functions are considered, M. Riesz gave examples
showing that the convexity result fails in the upper triangle and, yet, is
valid for the lower triangle. It follows immediately from the complex case,
however, that, given the hypotheses of Theorem~\ref{thm:c5-1-3}, there exists a
constant $k_{t}$ (not necessarily less than or equal to $k_{0}^{1-t}k_{1}^{t}$)
such that the operator $T$ is of type $(p_{t}, q_{t})$ with
$(p_{t}, q_{t})$-norm $\leq k_{t}$. In this connection, it is of interest to
remark that virtually all operators that arise in a natural way in Fourier
analysis are of type (or weak type) $(p, q)$ for some indices $p$ and $q$
satisfying $p \leq q$; thus, interpolation results corresponding to points of
the ``lower triangle'' suffice for our purpose.

\textbf{5.12.} When $p = 1$ and $1 < q \leq \infty$, then there is no norm
equivalent with $\|\cdot\|_{pq}^{*}$, and thus these spaces $L(p, q)$ cannot be
normed. Let us show this in the case $q = \infty$. On $E_{1}$ we consider the
sequence of functions $\{f_{k}\}$ defined by $f_{k}(x) = 1/|x - k|$. Then each
$f_{k} \in L(1, \infty)$, and $\|f_{k}\|_{1\infty}^{*}$ is independent of $k$.
Let $f = \sum_{k=1}^{N}f_{k}$. If $\|\cdot\|$ were a norm equivalent with
$\|\cdot\|_{1\infty}^{*}$, then we would have $\|f\| \leq C_{1}N$, for some
constant $C_{1}$. However, $f(x) \geq C_{2}\log N$ for $0 \leq x \leq N$, and
thus $\|f\|_{1\infty}^{*} \geq C_{2}N\log N$, giving a contradiction when
$N \to \infty$. The proof for $1 < q < \infty$ is similar.

\bibnoteshead

The Marcinkiewicz theorem was first announced in \cite{Marcinkiewicz2}, where
a sketch of the proof is given for the ``main diagonal''
$\{(\alpha, \beta) \in Q : \alpha = \beta\}$. Its proof for the ``lower
triangle,'' as well as many important applications was first published by
Zygmund \cite{Zygmund2}. The proof of the M. Riesz convexity theorem presented
in this chapter was discovered by Thorin \cite{Thorin1} and, independently, by
Tamarkin and Zygmund \cite{TamarkinZygmund1}. Interpolation of analytic
families of operators is due to Stein \cite{Stein2}. A related result, which
also makes use of Lemma~\ref{lem:c5-4-2} can be found in Hirschman
\cite{Hirschman1}. An extension of this type of interpolation to $H^{p}$ spaces
is found in Stein and G. Weiss \cite{SteinWeiss3}. $L(p, q)$ spaces were
introduced by Lorentz \cite{Lorentz1}. Calder\'on was the first to use the
function $m_{f}^{*}$ and to point out how it induces a norm on $L(p, q)$. For
the general theory of these spaces see Hunt \cite{Hunt1} and Oklander
\cite{Oklander1}. Our treatment of the results leading to the interpolation
Theorem~\ref{thm:c5-3-15} is adapted from Hunt \cite{Hunt1}. For related
results see Stein and G. Weiss \cite{SteinWeiss2}, Krein and Semenov
\cite{KreinSemenov1}, Calder\'on \cite{Calderon2} and Lions and Peetre
\cite{LionsPeetre1}. Two expository articles describing various general methods
of interpolation and how they are related to each other are Krein and Petunin
\cite{KreinPetunin1} and Magenes \cite{Magenes1}. Weak type inequalities for
the Hilbert transform go back to Kolmogorov \cite{Kolmogorov1}, Besicovitch
\cite{Besicovitch1} and Titchmarsh \cite{Titchmarsh1}.
